\documentclass[11pt]{article}

\usepackage[margin=1in]{geometry}
\usepackage[T1]{fontenc}
\usepackage{amsmath,amssymb,bm}
\usepackage{booktabs}
\usepackage{array}
\usepackage{longtable}
\usepackage{enumitem}
\usepackage[hidelinks]{hyperref}

% ---------------------------------------------------------------------------
% Notation macros (shared with verification_manual.tex where they overlap).
% Referential quantities carry the subscript R; the left Cauchy-Green tensor
% is bold B.
% ---------------------------------------------------------------------------
\newcommand{\ud}{\,\mathrm{d}}
\newcommand{\Grad}{\operatorname{Grad}}
\newcommand{\Div}{\operatorname{Div}}
\newcommand{\Dive}{\operatorname{Div}}
\newcommand{\grad}{\operatorname{grad}}
\newcommand{\dive}{\operatorname{div}}
\newcommand{\dev}{\operatorname{dev}}
\newcommand{\sym}{\operatorname{sym}}
\newcommand{\tr}{\operatorname{tr}}
\newcommand{\cof}{\operatorname{cof}}
\newcommand{\diag}{\operatorname{diag}}
\newcommand{\bx}{\bm{x}}
\newcommand{\bX}{\bm{X}}
\newcommand{\bu}{\bm{u}}
\newcommand{\bw}{\bm{w}}
\newcommand{\bb}{\bm{b}}
\newcommand{\be}{\bm{e}}
\newcommand{\bn}{\bm{n}}
\newcommand{\bN}{\bm{N}}
\newcommand{\bt}{\bm{t}}
\newcommand{\bT}{\bm{T}}
\newcommand{\bsig}{\bm{\sigma}}
\newcommand{\btau}{\bm{\tau}}
\newcommand{\beps}{\bm{\varepsilon}}
\newcommand{\bP}{\bm{P}}
\newcommand{\bS}{\bm{S}}
\newcommand{\bF}{\bm{F}}
\newcommand{\bFb}{\bar{\bm{F}}}
\newcommand{\bC}{\bm{C}}
\newcommand{\bB}{\bm{B}}
\newcommand{\bE}{\bm{E}}
\newcommand{\bI}{\bm{I}}
\newcommand{\bQ}{\bm{Q}}
\newcommand{\bH}{\bm{H}}
\newcommand{\bA}{\mathbb{A}}
\newcommand{\bphi}{\bm{\varphi}}
\newcommand{\bR}{\bm{R}}
\newcommand{\bK}{\bm{K}}
\newcommand{\bM}{\bm{M}}
\newcommand{\bBm}{\bm{\mathsf{B}}}
\newcommand{\Ione}{I_1}
\newcommand{\Itwo}{I_2}
\newcommand{\Ib}{\bar{I}_1}
\newcommand{\IIb}{\bar{I}_2}
\DeclareRobustCommand{\code}[1]{{\def\_{\textunderscore\allowbreak}\texttt{#1}}}
\DeclareRobustCommand{\file}[1]{{\def\_{\textunderscore\allowbreak}\texttt{#1}}}
\sloppy
\setlength{\emergencystretch}{2em}

\title{Continuum-Mechanics-MFEM\\[4pt]
\large Theory Manual of the Flux-Kernel Framework (\texttt{src/})}
\author{}
\date{September 2026}

\begin{document}
\maketitle

\begin{abstract}
\noindent
This manual describes the mathematical models and the numerical methods
implemented in the root \file{src/} library of the repository: finite-strain
solid mechanics in total Lagrangian form on a fixed reference mesh,
quasi-static or, when the input asks for it, with inertia, discretised with
continuous Galerkin finite elements on MFEM.  It
covers the kinematics and balance laws, the hyperelastic constitutive
library in four categories (compressible coupled, compressible uncoupled,
slightly compressible and incompressible), small-strain linear elasticity as
a special case of the same forms, the two-dimensional plane-strain and
plane-stress reductions, the displacement and the mixed
displacement--pressure formulations, boundary conditions and load
scheduling, the inertial term and its implicit time integration, the element
kernels and the dual-number tangents, the load and time stepper, damped
Newton, the algebraic multigrid and saddle-point linear solvers, and the
post-processing of quadrature-point quantities.  The verification and validation of the implementation against
closed-form and reference solutions is documented separately in the
verification manual (\file{doc/verification\_manual.tex}), which covers the
inputs and scripts under \file{apps/}.  The legacy applications under
\file{myapps/} are not part of this document.
\end{abstract}

\tableofcontents

%============================================================================
\section{Introduction}
%============================================================================

\subsection{Scope and architecture}

The framework is a multiphysics library on MFEM in which each physics is
meant to be expressed as fluxes and sources evaluated at quadrature points,
while the framework owns assembly, the nonlinear solver and the linear
solvers (the design is laid out in \file{doc/flux\_kernel\_architecture.html}
and the implementation plan in
\file{doc/hyperelasticity\_implementation\_plan.md}).  The first physics
is quasi-static nonlinear solid mechanics; the second, the transport of one
scalar unknown by diffusion, convection and reaction
(Section~\ref{sec:scalar}, plan \file{doc/scalar\_transport\_plan.md}),
is the first-order-in-time physics of the seam.  The library is organised
in five layers:

\begin{center}
\small
\begin{tabular}{@{}l p{11.5cm}@{}}
\toprule
Layer & Content \\
\midrule
\file{src/base/} & fixed-size tensors, forward-mode dual numbers, the YAML
  schema and its validation, the expression language for boundary data,
  mesh input, a named field registry, ParaView output and point probes \\
\file{src/kernels/} & the total Lagrangian element integrators
  (displacement and mixed $u$--$p$), the follower-pressure boundary
  face integrator, and the scalar flux kernel of the transport physics \\
\file{src/kernels/materials/} & the hyperelastic material functors and the
  small-strain linear elastic one, the kinematics trait, the plane-stress
  adapter, the dual-number tangent, the material variant and the YAML
  factory; the affine laws of the scalar transport \\
\file{src/physics/} & the solid mechanics modules (residual, Jacobian,
  fields), the load set shared by both formulations, the
  quadrature-point post-processing, and the scalar transport module with
  the conditions of a scalar unknown \\
\file{src/solvers/} & the quasi-static load stepper, damped Newton, the
  Krylov--AMG linear solver, the augmented-Lagrangian saddle-point solver
  and the sparse direct solver \\
\bottomrule
\end{tabular}
\end{center}

Executables under \file{apps/} only parse the input and wire library
objects together; they contain no physics, assembly or solver code.  The
executables \file{apps/solid\_mechanics.cpp} and \file{apps/scalar\_transport.cpp}
and their inputs are described in the verification manual.

\subsection{What is implemented}

\begin{itemize}[nosep]
  \item Balance of linear momentum in total Lagrangian form: the mesh never
    moves, every integral is over the reference configuration.  By default
    the analysis is quasi-static: loads and prescribed displacements are
    functions of a pseudo-time $t \in [0, 1]$ and inertia is absent.
  \item Elastodynamics on request (Section~\ref{sec:dynamics}): the inertial
    term $\int_{\Omega_R}\rho_R\ddot\bu\cdot\bw\ud V$ with a constant,
    consistent mass matrix and the reference density by element attribute,
    integrated in physical time by the generalized-$\alpha$ family in
    displacement form (Newmark, HHT-$\alpha$, generalized-$\alpha$), for
    every material and both formulations.  It is a decorator over the
    quasi-static problem, not a second set of kernels.
  \item Hyperelasticity in four categories (Section~\ref{sec:classification}):
    compressible coupled models (neo-Hookean of Bonet--Wood form, Saint
    Venant--Kirchhoff); compressible uncoupled models
    $\Psi_{\mathrm{iso}}(\bFb) + \kappa\,u(J)$ with a selectable volumetric
    law $u(J)$, quadratic $\tfrac12(J-1)^2$ by default (isochoric
    neo-Hookean, Mooney--Rivlin, Yeoh, Gent, Arruda--Boyce, Ogden); the
    same six as slightly compressible materials with the pressure as an
    independent field; and the same six as incompressible materials with
    $\kappa = \infty$ and a Lagrange multiplier.  Materials are stateless
    and may vary by element attribute.
  \item Small-strain (geometrically linear) isotropic elasticity,
    \code{linear\_elastic}, in both formulations, including $\nu = \tfrac12$
    (Section~\ref{sec:smallstrain}).  It is a material of the same kernels,
    with the kinematics as a compile-time trait; the problem is linear and
    is solved by one Newton step.
  \item Finite viscoelasticity (Section~\ref{sec:viscoelasticity}): Maxwell
    branches over any decoupled model, with the viscous right Cauchy--Green
    tensor of each branch as a quadrature-point internal variable, integrated
    by the reference's implicit update, in both formulations and with
    inertia; the first history-dependent material of the library, with a
    quasi-static analysis in physical time (the \code{time} block).
  \item Finite thermoelasticity (Section~\ref{sec:thermoelasticity}): a
    decoupled model with the temperature as a second state variable
    (entropic shear modulus, thermal expansion through the volumetric law,
    Fourier conduction) and the heat equation with the thermoelastic
    heating, solved as a coupled displacement--pressure--temperature problem
    on Taylor--Hood spaces in physical time, with prescribed temperatures
    and heat fluxes per current area as boundary conditions.
  \item Scalar transport (Section~\ref{sec:scalar}): the
    convection--diffusion--reaction equation of one unknown on a scalar
    $H^1$ space, with a capacity and a conductivity affine in the unknown,
    a velocity and a source as expressions, the convection in
    non-conservative or conservative form, implicit Euler in physical time
    or a steady solve, prescribed values and inward fluxes with the flows
    through the Dirichlet entries, the errors against an exact expression,
    and one solve for a linear problem; the second physics of the seam,
    with its own executable.
  \item Two formulations: a pure displacement formulation with the
    assembled tangent, and a mixed displacement--pressure formulation on
    Taylor--Hood pairs for near- and fully incompressible materials; the
    coupled formulation extends the mixed one by the temperature.
  \item Two-dimensional problems as plane strain (default), plane stress
    (displacement formulation; the thickness stretch is eliminated
    pointwise at the material level) or axisymmetric (Section~\ref{sec:axisymmetric}:
    the meridian section of a solid of revolution, the hoop stretch in the
    kernels and the weight $2\pi r$ on every integral, either formulation).
  \item Boundary conditions: prescribed displacements on all or some
    components of faces or of single nodes, nominal tractions, dead normal
    pressures, follower pressures per unit current area, a body force per
    unit mass, a rigid-sphere penalty contact, and for the coupled problem
    prescribed temperatures and inward heat fluxes, each with its own load
    schedule and data given as expressions of the reference coordinates and
    the pseudo-time.
  \item Consistent tangents from forward-mode automatic differentiation of
    the same code that evaluates the stress (Ogden and the compressible
    plane-stress adapter use analytic or implicit derivatives where
    eigenvectors or a root are not differentiable by duals alone).
  \item Solvers: load stepping, and time stepping, with bisection on
    failure, damped Newton with Armijo backtracking, GMRES or CG with
    BoomerAMG, and FGMRES with a block upper-triangular preconditioner for
    the $u$--$p$ saddle point, whose Schur complement approximation gains an
    inertial part in a dynamic analysis.
  \item Post-processing of Cauchy and first Piola--Kirchhoff stress, the
    deformation gradient, the strain, the volume ratio, von Mises stress, energy density
    and thickness stretch at quadrature points, presented as nodal fields,
    element averages or point clouds; probes; reactions of the Dirichlet
    entries; in a dynamic analysis the nodal velocity and acceleration, and
    the kinetic and internal energy with the external work.
\end{itemize}

Not implemented: explicit time integration, physical damping other than
the Maxwell branches, modal analysis, inelasticity other than the viscous
flow of those branches (plasticity, damage), inertia with a temperature
field, convection or radiation boundary conditions, volumetric heat
sources, temperature-dependent conductivity or heat capacity, torsion in
the axisymmetric description, contact other than the rigid-sphere penalty,
point loads, multi-point or periodic constraints, discontinuous Galerkin or
hybridised discretisations, partial assembly, and device execution; for the
scalar transport, stabilisation (SUPG), anisotropic or tabulated
coefficients, Robin conditions, higher-order time integration, regions, and
the ALE description.  Section~\ref{sec:seam} lists what of the present code
is specific to each physics and what the two share.

\subsection{Conventions and companion notes}

Referential (Lagrangian) quantities carry the subscript $R$; the current
configuration is $\Omega_t$.  The left Cauchy--Green tensor is $\bB = \bF\bF^T$,
the right one $\bC = \bF^T\bF$.  The pressure unknown of the mixed formulation
and the \code{pressure} output field are the tensile-positive mean stress
$p = \tfrac13\tr\bsig$.  In two dimensions every $3\times3$ tensor is built
from the in-plane block with $F_{33} = 1$ (plane strain), $F_{33} = \lambda_3$
(plane stress) or $F_{33} = 1 + u_r/r$ (axisymmetric).

The verification manual (\file{doc/verification\_manual.tex}) documents
every input and test; its appendix collects the closed-form homogeneous
states of the incompressible models that serve as reference solutions.

%============================================================================
\section{Kinematics}
%============================================================================

Let $\Omega_R \subset \mathbb{R}^d$, $d \in \{2, 3\}$, be the reference
configuration with boundary $\partial\Omega_R$ and outward unit normal
$\bN$.  The motion $\bx = \bphi(\bX, t) = \bX + \bu(\bX, t)$ carries the
material point $\bX$ to $\bx$; $t$ is a pseudo-time that parametrises the
load path (Section~\ref{sec:loads}).  The deformation gradient, its
determinant and the strain measures are
\begin{gather}
  \bF = \Grad\bphi = \bI + \bH, \qquad \bH = \Grad\bu, \qquad
  J = \det\bF > 0, \\
  \bC = \bF^T\bF, \qquad \bB = \bF\bF^T, \qquad
  \bE = \tfrac12(\bC - \bI),
\end{gather}
with $\Grad = \partial/\partial\bX$.  The principal invariants of $\bC$
(equivalently of $\bB$) are
\begin{equation}
  \Ione = \tr\bC, \qquad
  \Itwo = \tfrac12\bigl[(\tr\bC)^2 - \tr(\bC^2)\bigr], \qquad
  I_3 = \det\bC = J^2,
\end{equation}
and in terms of the principal stretches $\lambda_a$ (square roots of the
eigenvalues $c_a$ of $\bC$) $\Ione = \sum_a\lambda_a^2$,
$\Itwo = \sum_{a<b}\lambda_a^2\lambda_b^2$, $J = \lambda_1\lambda_2\lambda_3$.
The isochoric (volume-preserving) part of the deformation and the modified
invariants are
\begin{equation}
  \bFb = J^{-1/3}\bF, \qquad
  \Ib = J^{-2/3}\Ione, \qquad
  \IIb = J^{-4/3}\Itwo, \qquad
  \bar\lambda_a = J^{-1/3}\lambda_a .
  \label{eq:isochoric}
\end{equation}

The stress measures are the Cauchy stress $\bsig$ (force per unit current
area), the first Piola--Kirchhoff (nominal) stress $\bP$ (force per unit
reference area) and the second Piola--Kirchhoff stress $\bS$, related by
\begin{equation}
  \bP = J\bsig\bF^{-T}, \qquad
  \bS = \bF^{-1}\bP, \qquad
  \bsig = J^{-1}\bP\bF^T .
  \label{eq:piola}
\end{equation}
$\bsig$ and $\bS$ are symmetric; $\bP$ is not, but satisfies
$\bP\bF^T = \bF\bP^T$.  Nanson's formula relates area elements,
\begin{equation}
  \bn\ud a = J\bF^{-T}\bN\ud A = \cof\bF\,\bN\ud A,
  \label{eq:nanson}
\end{equation}
and is what turns a traction per unit current area into a nominal
traction.  The reference density is $\rho_R$; mass conservation gives
$\rho_R = J\rho$.

In two dimensions the kernels form the $2\times2$ displacement gradient
$\bH_{2D}$ and embed it as
\begin{equation}
  \bF = \begin{pmatrix} \bI_2 + \bH_{2D} & \bm 0 \\ \bm 0^T & F_{33} \end{pmatrix},
  \qquad F_{33} = 1 \ \text{(plane strain)}, \qquad
  F_{33} = \lambda_3 \ \text{(plane stress, Section~\ref{sec:planestress})},
  \label{eq:embed}
\end{equation}
so that one three-dimensional constitutive code path serves every
dimension.

%============================================================================
\section{Balance of momentum and the displacement weak form}
\label{sec:momentum}
\label{sec:weak}
%============================================================================

\subsection{Strong form}

Quasi-static balance of linear momentum on the reference configuration
reads
\begin{subequations}
\label{eq:strong}
\begin{alignat}{2}
  \Dive\bP + \rho_R\bb &= \bm 0 &\quad& \text{in } \Omega_R, \\
  \bu &= \bar\bu && \text{on } \Gamma_R^D, \\
  \bP\bN &= \bar\bT && \text{on } \Gamma_R^N,
\end{alignat}
\end{subequations}
where $\bb$ is the body force per unit mass, $\bar\bu$ the prescribed
displacement (possibly of a subset of components) and $\bar\bT$ the
prescribed nominal traction.  Balance of angular momentum is
$\bP\bF^T = \bF\bP^T$ and is satisfied identically by hyperelastic stresses.
The system is closed by a hyperelastic constitutive law
\begin{equation}
  \bP = \frac{\partial\Psi}{\partial\bF}(\bF),
  \label{eq:hyper}
\end{equation}
with $\Psi$ the stored energy per unit reference volume
(Section~\ref{sec:materials}).

\subsection{Weak form}

With the trial and test spaces on the fixed reference domain,
\begin{equation}
  \mathcal S = \{\bu \in [H^1(\Omega_R)]^d : \bu = \bar\bu \text{ on } \Gamma_R^D\},
  \qquad
  \mathcal V = \{\bw \in [H^1(\Omega_R)]^d : \bw = \bm 0 \text{ on } \Gamma_R^D\},
\end{equation}
the total Lagrangian weak form is: find $\bu \in \mathcal S$ such that for
all $\bw \in \mathcal V$
\begin{equation}
  \boxed{\;
  R(\bu; \bw) \;=\;
  \int_{\Omega_R}\bP(\bF):\Grad\bw\ud V
  \;-\;\int_{\Omega_R}\rho_R\,\bb\cdot\bw\ud V
  \;-\;\int_{\Gamma_R^N}\bar\bT\cdot\bw\ud A
  \;+\; R_{\text{fol}}(\bu; \bw)
  \;=\; 0 ,\;}
  \label{eq:weak}
\end{equation}
where $R_{\text{fol}}$ is the contribution of follower pressures
(Section~\ref{sec:follower}), zero when there are none.  The sign convention
throughout the code is \emph{residual = internal minus external}.  The form
is nonlinear in $\bu$ through $\bP(\bF)$ and through the geometry in $\bF$.

\paragraph{Equivalent forms.}
With $\bP = \bF\bS$ and the symmetry of the second Piola--Kirchhoff stress,
$\bP:\Grad\bw = \bS:\delta\bE[\bw]$, where
$\delta\bE[\bw] = \tfrac12(\bF^T\Grad\bw + \Grad^T\bw\,\bF)$ is the variation
of the Green--Lagrange strain, so the internal virtual work may be written
$\int_{\Omega_R}\bS(\bE):\delta\bE[\bw]\ud V$, the form commonly used with
$(\bS, \bE)$ constitutive pairs such as Saint Venant--Kirchhoff; its
linearisation splits into the familiar material and geometric stiffness,
which the nominal tangent $\bA$ of \eqref{eq:tangent} contains together.
Pushing forward with $\bP:\Grad\bw = J\bsig:\grad\bw$ and
$\ud v = J\ud V$ gives the spatial statement on the current configuration
$\Omega_t$,
\begin{equation}
  \int_{\Omega_t}\bsig:\grad\bw\ud v
  = \int_{\Omega_t}\rho\,\bb\cdot\bw\ud v + \int_{\Gamma_t^N}\bar\bt\cdot\bw\ud a ,
  \qquad
  \dive\bsig + \rho\bb = \bm 0 \ \text{in } \Omega_t, \quad
  \bsig\bn = \bar\bt \ \text{on } \Gamma_t^N, \quad \bsig = \bsig^T,
\end{equation}
with $\rho = \rho_R/J$ and $\bar\bt$ the traction per unit current area: the
two statements are identical and differ only in the configuration on which
the integrals are evaluated.  The code integrates on $\Omega_R$ throughout.
With inertia the residual gains $\int_{\Omega_R}\rho_R\ddot\bu\cdot\bw\ud V$
(spatially $\rho\dot{\bm v}$ with the material acceleration
$\dot{\bm v} = \partial_t\bm v + (\grad\bm v)\bm v$) and needs initial data
$\bu_0$, $\bm v_0$: Section~\ref{sec:dynamics}.

\subsection{Linearisation}

Newton's method needs the directional derivative of \eqref{eq:weak}.  For
the internal term,
\begin{equation}
  D_{\Delta\bu}\int_{\Omega_R}\bP:\Grad\bw\ud V
  = \int_{\Omega_R}\Grad\bw:\bA:\Grad\Delta\bu\ud V,
  \qquad
  \bA = \frac{\partial\bP}{\partial\bF}
      = \frac{\partial^2\Psi}{\partial\bF\,\partial\bF},
  \quad A_{ijkl} = \frac{\partial P_{ij}}{\partial F_{kl}} .
  \label{eq:tangent}
\end{equation}
$\bA$ is the fourth-order material tangent in nominal form.  For a
hyperelastic material it has the major symmetry $A_{ijkl} = A_{klij}$, so
the assembled stiffness is symmetric; the follower pressure adds a
non-symmetric part (Section~\ref{sec:follower}).  The code never writes
$\bA$ by hand: it is produced by differentiating the stress evaluation with
dual numbers (Section~\ref{sec:ad}).

%============================================================================
\section{Constitutive models}
\label{sec:materials}
%============================================================================

Every material in the library is a stateless value type whose parameters
are public members and whose stress function is a template over the scalar
type,
\begin{equation}
  \bP = \mathrm{PK1}(\bF): \ \text{tensor}\langle T, 3, 3\rangle \to
  \text{tensor}\langle T, 3, 3\rangle, \qquad T \in \{\code{double}, \code{dual}\},
\end{equation}
optionally with an energy $\Psi(\bF)$.  Models for the mixed formulation
expose the split $\bP_{\mathrm{iso}}$, $\Psi_{\mathrm{iso}}$, the
volumetric pressure $U'(J)$, the bulk modulus $\kappa$ (possibly infinite)
and a small-strain shear modulus used to scale the saddle-point
preconditioner.

\subsection{Classification by the treatment of volume change}
\label{sec:classification}

All models are isotropic and hyperelastic; what distinguishes them is how
the volume change $J$ enters the energy and how the resulting pressure is
handled numerically.  Four categories are used in this manual:

\begin{center}
\small
\renewcommand{\arraystretch}{1.2}
\begin{tabular}{@{}>{\raggedright\arraybackslash}p{2.5cm} >{\raggedright\arraybackslash}p{4.6cm} >{\raggedright\arraybackslash}p{3.2cm} >{\raggedright\arraybackslash}p{4.4cm}@{}}
\toprule
Category & Volumetric treatment & Models & Formulation and input \\
\midrule
1. Compressible, coupled &
  one energy $\Psi(\bF)$; $J$ enters through $\ln J$ or $\tr\bE$ terms mixed with the shear response &
  \code{neo\_hookean}, \code{st\_venant\_kirchhoff}, \code{gent\_compressible\_summit} &
  displacement; \code{E}, \code{nu} with $\nu < \tfrac12$ \\
2. Compressible, uncoupled &
  $\Psi_{\mathrm{iso}}(\bFb) + \kappa\,u(J)$, $u$ one of four laws (quadratic $\tfrac12(J-1)^2$ by default); $\kappa/\mu$ of order one to ten &
  \code{iso\_neo\_hookean}, \code{mooney\_rivlin}, \code{yeoh}, \code{gent}, \code{arruda\_boyce}, \code{ogden} &
  displacement (mixed also allowed); \code{kappa} or \code{nu} \\
3. Slightly compressible &
  the same split with $\kappa \gg \mu$; the pressure $p = U'(J)$ becomes an independent field &
  the same six &
  mixed (the displacement formulation locks); \code{kappa} or \code{nu} close to $\tfrac12$ \\
4. Incompressible &
  $J = 1$ enforced by a Lagrange multiplier $p$; $\kappa = \infty$ &
  the same six &
  mixed, or plane stress in 2D; \code{incompressible: true} or \code{nu: 0.5} \\
\bottomrule
\end{tabular}
\end{center}

Categories 2, 3 and 4 share one code path: the six decoupled materials
provide $\Psi_{\mathrm{iso}}$ and $U'(J)$, and the category is set by the
value of $\kappa$ and by the formulation chosen in the input.  The border
between 2 and 3 is one of degree, not of kind: the displacement
formulation is adequate up to $\kappa/\mu$ of order ten ($\nu \approx 0.45$,
where the tests confirm that both formulations agree under refinement),
while at $\nu = 0.4999$ ($\kappa/\mu \approx 5000$) the mixed formulation
is required in practice.  Category 1 is a different code path with no
split; those two models cannot be made incompressible.

\subsection{Category 1: compressible, coupled}
\label{sec:coupled}

\paragraph{Neo-Hookean (Bonet--Wood form).}
\begin{equation}
  \Psi = \frac{\mu}{2}(\Ione - 3) - \mu\ln J + \frac{\lambda}{2}(\ln J)^2,
  \qquad
  \bP = \mu(\bF - \bF^{-T}) + \lambda\ln J\,\bF^{-T} .
  \label{eq:nh}
\end{equation}
This is the compressible extension of the neo-Hookean energy given by
Bonet and Wood~\cite{bonetwood} and used by the FEniCS and FEniCSx
hyperelasticity demos and by FEBio (the source comment also names
Simo--Hughes; the exact expression is not confirmed there).  The
$-\mu\ln J$ term cancels the stress of the $\mu\bF$ term at $\bF = \bI$, so
the reference state is stress free; the small-strain limit is linear
elasticity with the Lam\'e constants $(\mu, \lambda)$; and both logarithmic
terms diverge as $J \to 0$, so compression to zero volume costs infinite
energy.  The shear term uses the full $\Ione$, not $\Ib$, so shear and
volume are coupled and the volumetric response is $\lambda\ln J$ rather
than a separate law $U(J)$.  Other coupled compressible neo-Hookean forms
in the literature replace $\tfrac{\lambda}{2}(\ln J)^2$ by
$\tfrac{\lambda}{2}(J-1)^2$ or by the polyconvex Ciarlet--Geymonat
combination $\tfrac{\Lambda}{4}(J^2 - 1) - \tfrac{\Lambda}{2}\ln J$
\cite{pencegou}; none of these is implemented.

\paragraph{Saint Venant--Kirchhoff.}
\begin{equation}
  \Psi = \frac{\lambda}{2}(\tr\bE)^2 + \mu\,\bE:\bE,
  \qquad
  \bS = \lambda\tr\bE\,\bI + 2\mu\bE,
  \qquad
  \bP = \bF\bS .
  \label{eq:svk}
\end{equation}
This is linear elasticity in the Green--Lagrange strain; it is not
polyconvex, its energy stays finite as $J \to 0$ and it softens in
compression, so it is used chiefly for manufactured-solution tests.

\paragraph{Compressible Gent, SUMMIT form (\code{gent\_compressible\_summit}).}
\begin{equation}
  \Psi = -\frac{\mu}{2}\Bigl(J_m\ln\Bigl(1 - \frac{\Ione - 3}{J_m}\Bigr) + 2\ln J\Bigr)
         + \frac{\kappa}{2}A^4,
  \qquad
  A = \frac{J^2 - 1}{2} - \ln J,
  \label{eq:gentsummit}
\end{equation}
\begin{equation}
  \bP = \frac{\mu J_m}{J_m - \Ione + 3}\,\bF
        + \bigl(2\kappa A^3(J^2 - 1) - \mu\bigr)\bF^{-T},
  \qquad \Ione - 3 < J_m .
\end{equation}
This is the \code{gent-compressible} material of the SUMMIT code
(\code{summit::GentCompressibleHyperelastic}), transcribed so that results
can be compared between the two codes: the Gent limited-extensibility term
on the full $\Ione$ (not $\Ib$, unlike the decoupled \code{gent} of
Section~\ref{sec:split}), the $-\mu\ln J$ term that makes the reference
state stress free, and a quartic volumetric penalty.  For $J_m \to \infty$
the chain term reduces to the neo-Hookean $\tfrac{\mu}{2}(\Ione - 3)$.
Since $A = (J-1)^2 + O((J-1)^3)$, the penalty is $O((J-1)^8)$ and $\kappa$
does \emph{not} enter the small-strain response: the linearised moduli are
$\mu$ and $\lambda = 2\mu/J_m$ (the latter from the chain term, because
$\Ione - 3 \approx 2\tr\beps$), that is $\nu = 1/(J_m + 2)$, whatever
$\kappa$ is.  $\kappa$ is therefore a penalty parameter that stiffens the
volume change at finite strain only, not a bulk modulus; SUMMIT reports
$E$ and $\nu$ computed from $\kappa$ and $\mu$ as if it were one (they enter
only its wave-speed estimate).  The input keys are \code{mu}, \code{kappa}
and \code{Jm}; the tangent comes from automatic differentiation and agrees
with SUMMIT's closed-form moduli to round-off (\file{test\_materials}).
Displacement formulation only.

The first two take $(E, \nu)$ from the input and convert with
$\mu = E/(2(1+\nu))$, $\lambda = E\nu/((1+\nu)(1-2\nu))$; the derived bulk
modulus $\kappa = \lambda + \tfrac23\mu$ is reported but does not enter the
stress.  None of the three models separates the pressure, so none can be
used in the mixed formulation or in the incompressible limit.

\subsection{Category 2: compressible, uncoupled}
\label{sec:split}

The remaining six models are decoupled,
\begin{equation}
  \Psi(\bF) = \Psi_{\mathrm{iso}}(\bFb) + U(J),
  \qquad
  U(J) = \kappa\,u(J),
  \qquad
  u(J) = \tfrac12(J - 1)^2 \ \text{by default},
  \label{eq:split}
\end{equation}
where $\Psi_{\mathrm{iso}}$ depends on $\bF$ only through the isochoric
part \eqref{eq:isochoric} and $u(J)$ is one of the volumetric laws of the
paragraph below, normalised so that $u''(1) = 1$.  Because $\Psi_{\mathrm{iso}}$ is invariant
under $\bF \to \alpha\bF$, Euler's theorem gives
$\bP_{\mathrm{iso}}:\bF = 0$, i.e.\ the isochoric Cauchy stress
$J^{-1}\bP_{\mathrm{iso}}\bF^T$ is deviatoric for every $\bF$, not only at
$J = 1$.  In the displacement formulation the stress is
\begin{equation}
  \bP = \bP_{\mathrm{iso}}(\bF) + U'(J)\,J\bF^{-T},
  \qquad
  \bsig = J^{-1}\bP_{\mathrm{iso}}\bF^T + U'(J)\,\bI ,
  \qquad
  U'(J) = \kappa\,u'(J) = \kappa(J-1) \ \text{by default},
  \label{eq:Pdisp}
\end{equation}
so $p = U'(J)$ is exactly the mean stress $\tfrac13\tr\bsig$.  At
moderate $\kappa/\mu$ this is an ordinary compressible material and the
pure displacement discretisation is well behaved.

\paragraph{Volumetric law.}
The volumetric law is selectable per material (\code{material.volumetric},
\file{materials/volumetric.hpp}).  Writing $U(J) = \kappa\,u(J)$, the four
laws are
\begin{center}
\small
\renewcommand{\arraystretch}{1.4}
\begin{tabular}{@{}l l l l l@{}}
\toprule
Key & $u(J)$ & $u'(J)$ & $u''(J)$ & Notes \\
\midrule
\code{quadratic} & $\tfrac12(J-1)^2$ & $J - 1$ & $1$ & default; $U$ finite as $J \to 0$ \\
\code{simo\_taylor} & $\tfrac14(J^2 - 1 - 2\ln J)$ & $\tfrac12(J - J^{-1})$ & $\tfrac12(1 + J^{-2})$ & Simo--Taylor; convex, $U \to \infty$ both ways \\
\code{logarithmic} & $\tfrac12(\ln J)^2$ & $\ln J/J$ & $(1 - \ln J)/J^2$ & Anand's FEniCSx codes; $u'$ peaks at $J = e$ \\
\code{j\_log\_j} & $J\ln J - J + 1$ & $\ln J$ & $1/J$ & polyconvex; $u'$ monotone \\
\bottomrule
\end{tabular}
\end{center}
The logarithmic law softens in dilatation: its tangent bulk modulus
$\kappa(1 - \ln J)/J^2$ falls to $0.18\,\kappa$ at $J = 1.67$ and to zero at
$J = e$, where the pressure peaks at $\kappa/e$.  Its Newton iterates must
therefore stay near $J = 1$ (run it with the tangent predictor of
Section~\ref{sec:solvers}), and it can fail at stress singularities where
$J$ grows locally, whatever the increment size; the monotone laws do not
have this limitation.
Every law satisfies $u(1) = u'(1) = 0$ and $u''(1) = 1$, so $\kappa$ is the
small-strain bulk modulus for all of them and the input semantics of
\code{kappa} and \code{nu} do not change; the laws agree to first order in
the strain and differ at finite strain, chiefly in strong compression and
in the pressure of nearly incompressible problems (Doll and
Schweizerhof~\cite{doll} survey the family).  The displacement formulation
needs only $U'(J) = \kappa u'(J)$; the mixed formulation uses the
normalised derivatives $u'$ and $u''$
(Sections~\ref{sec:nearly} and~\ref{sec:mixed}).  Regions inherit the law
of the base material.  The law is irrelevant at $\kappa = \infty$
(category 4) and is not a key of the coupled models.

\paragraph{Isochoric energies.}
The six isochoric energies below are shared by categories 2, 3 and 4.

\subsubsection{Models depending on $\Ib$ alone}

For $\Psi_{\mathrm{iso}} = f(\Ib)$,
\begin{equation}
  \bP_{\mathrm{iso}} = f'(\Ib)\,\frac{\partial\Ib}{\partial\bF},
  \qquad
  \frac{\partial\Ib}{\partial\bF}
  = J^{-2/3}\Bigl(2\bF - \tfrac23\Ione\,\bF^{-T}\Bigr),
  \label{eq:dI1bar}
\end{equation}
so that a model only has to provide $f'$ (\file{isochoric.hpp}).  The
library contains
\begin{center}
\small
\renewcommand{\arraystretch}{1.6}
\begin{tabular}{@{}l l l l@{}}
\toprule
Model & $\Psi_{\mathrm{iso}}$ & $f'(\Ib) = \partial\Psi_{\mathrm{iso}}/\partial\Ib$ & small-strain $\mu$ \\
\midrule
Isochoric neo-Hookean & $\frac{\mu}{2}(\Ib - 3)$ & $\frac{\mu}{2}$ & $\mu$ \\
Yeoh & $\sum_{i=1}^3 c_{i0}(\Ib - 3)^i$ & $c_{10} + 2c_{20}(\Ib-3) + 3c_{30}(\Ib-3)^2$ & $2c_{10}$ \\
Gent & $-\frac{\mu J_m}{2}\ln\bigl(1 - \frac{\Ib - 3}{J_m}\bigr)$ & $\frac{\mu}{2}\frac{J_m}{J_m - (\Ib - 3)}$ & $\mu$ \\
Arruda--Boyce (Pad\'e) & $\frac{\mu}{6}(\Ib - 3) - \mu N\ln\frac{3N - \Ib}{3N - 3}$ & $\frac{\mu}{6}\,\frac{9N - \Ib}{3N - \Ib}$ & $\mu\,\frac{3N - 1}{3N - 3}$ \\
Arruda--Boyce (series) & eq.~\eqref{eq:abseries} & eq.~\eqref{eq:abseries} & $\mu_0$, eq.~\eqref{eq:abmu0} \\
\bottomrule
\end{tabular}
\end{center}
\emph{Isochoric neo-Hookean} (\code{iso\_neo\_hookean}): the Gaussian-chain
model; the closed form is
$\bP_{\mathrm{iso}} = \mu J^{-2/3}(\bF - \tfrac13\Ione\bF^{-T})$.
\emph{Yeoh} (\code{yeoh}): a cubic in $\Ib - 3$; typically $c_{20} < 0$,
$c_{30} > 0$, reproducing the softening-then-stiffening response of filled
rubbers.  \emph{Gent} (\code{gent}): limited chain extensibility; the
energy is defined only for $\Ib - 3 < J_m$, beyond the locking stretch it
evaluates to NaN, which the Newton line search rejects like any non-finite
residual, and $J_m \to \infty$ recovers the neo-Hookean model.
\emph{Arruda--Boyce} (\code{arruda\_boyce}): the eight-chain model,
\begin{equation}
  f'(\Ib) = \mu\,\frac{\mathcal{L}^{-1}(z)}{6z}, \qquad z = \sqrt{\frac{\Ib}{3N}},
  \label{eq:abexact}
\end{equation}
with the inverse Langevin function $\mathcal{L}^{-1}$ of the relative
chain stretch $z$; $\mu$ is the chain-density parameter $nkT$ and $N$ the
number of rigid links per chain (locking stretch $\lambda_L = \sqrt N$).
$\mathcal{L}^{-1}$ has no closed form, and the key
\code{inverse\_langevin} selects its approximation.  The default,
\code{pade}, is Cohen's Pad\'e approximation~\cite{cohen}
$\mathcal{L}^{-1}(z) \approx z(3 - z^2)/(1 - z^2)$, which keeps the
singularity at $z = 1$:
\begin{equation}
  f'(\Ib) = \frac{\mu}{6}\,\frac{9N - \Ib}{3N - \Ib}, \qquad
  \Psi_{\mathrm{iso}} = \frac{\mu}{6}(\Ib - 3) - \mu N\ln\frac{3N - \Ib}{3N - 3}, \qquad
  \mu_0 = 2f'(3) = \mu\,\frac{3N - 1}{3N - 3}.
  \label{eq:abpade}
\end{equation}
The energy is the integral of $f'$, so stress, tangent and the energy of
the line search are consistent; it needs $N > 1$ and is defined for
$\Ib < 3N$, beyond which it evaluates to NaN and the line search rejects
the step, as for the Gent model.  The Pad\'e form is within $5\%$ of the
exact $\mathcal{L}^{-1}$ for all $z$ and exact at $z \to 1$.  With
\code{inverse\_langevin: series} the model is the five-term expansion of
\eqref{eq:abexact} about $z = 0$,
\begin{equation}
  \Psi_{\mathrm{iso}} = \mu\sum_{i=1}^5\frac{C_i}{N^{i-1}}(\Ib^{\,i} - 3^i), \qquad
  f'(\Ib) = \mu\sum_{i=1}^5\frac{i\,C_i}{N^{i-1}}\Ib^{\,i-1}, \qquad
  C = \{\tfrac12, \tfrac1{20}, \tfrac{11}{1050}, \tfrac{19}{7000}, \tfrac{519}{673750}\},
  \label{eq:abseries}
\end{equation}
with the small-strain shear modulus
\begin{equation}
  \mu_0 = \mu\Bigl(1 + \frac{3}{5N} + \frac{99}{175N^2} + \frac{513}{875N^3} + \frac{42039}{67375N^4}\Bigr).
  \label{eq:abmu0}
\end{equation}
The series has lost the singularity and is softer near the locking
stretch.  The two differ in $\mu_0$ at order $1/N$, because the Pad\'e form
gives $3z + 2z^3$ at small $z$ against the exact $3z + \tfrac95 z^3$.

\subsubsection{Mooney--Rivlin}

\begin{gather}
  \Psi_{\mathrm{iso}} = c_1(\Ib - 3) + c_2(\IIb - 3), \qquad
  \Itwo = \tfrac12(\Ione^2 - \bC:\bC), \\
  \bP_{\mathrm{iso}}
  = c_1 J^{-2/3}\Bigl(2\bF - \tfrac23\Ione\bF^{-T}\Bigr)
  + c_2 J^{-4/3}\Bigl(2\Ione\bF - 2\bF\bC - \tfrac43\Itwo\bF^{-T}\Bigr),
  \qquad \mu = 2(c_1 + c_2).
\end{gather}
The only model in the library that depends on the second invariant.  With
$c_2 = 0$ it is the isochoric neo-Hookean model.

\subsubsection{Ogden}
\label{sec:ogden}

The Ogden energy is a function of the isochoric principal stretches,
\begin{equation}
  \Psi_{\mathrm{iso}} = \sum_{r=1}^{M}\frac{\mu_r}{\alpha_r}
    \bigl(\bar\lambda_1^{\alpha_r} + \bar\lambda_2^{\alpha_r} + \bar\lambda_3^{\alpha_r} - 3\bigr),
  \qquad M \le 6, \quad \mu_r\alpha_r > 0, \qquad
  \mu = \tfrac12\sum_r\mu_r\alpha_r .
\end{equation}
It is evaluated spectrally.  With $\bC = \sum_a c_a\bN_a\otimes\bN_a$
(cyclic Jacobi eigen-decomposition, \file{spectral.hpp}) define
\begin{equation}
  \beta_a = \sum_r\mu_r\bar\lambda_a^{\alpha_r}, \qquad
  \tau_a = \beta_a - \tfrac13(\beta_1 + \beta_2 + \beta_3), \qquad
  \bS_{\mathrm{iso}} = \sum_a\frac{\tau_a}{c_a}\bN_a\otimes\bN_a, \qquad
  \bP_{\mathrm{iso}} = \bF\bS_{\mathrm{iso}} ,
  \label{eq:ogden}
\end{equation}
where $\tau_a$ are the isochoric principal Kirchhoff stresses (the
$\tfrac13\sum\beta$ term is the deviatoric projection that
$\bar\lambda_a = J^{-1/3}\lambda_a$ produces).  Eigenvectors are not
differentiable at coincident eigenvalues, so the tangent is not obtained by
running dual numbers through the eigen-solver.  Instead the dual overload
of $\bP_{\mathrm{iso}}$ forms the directional derivative along $\ud\bF$
analytically (Ogden 1984; Simo and Taylor 1991): with
$\ud\bC = \ud\bF^T\bF + \bF^T\ud\bF$ expressed in the eigenbasis,
$\ud\hat{\bC} = \bN^T\ud\bC\,\bN$, the derivative of $\bS_{\mathrm{iso}}$
in that basis has diagonal entries $\sum_b (\partial s_a/\partial c_b)\ud\hat C_{bb}$
with $s_a = \tau_a/c_a$, and off-diagonal entries
\begin{equation}
  \ud\hat S_{ab} = \gamma_{ab}\,\ud\hat C_{ab}, \qquad
  \gamma_{ab} = \begin{cases}
    \dfrac{s_a - s_b}{c_a - c_b}, & |c_a - c_b| > 10^{-8}\max_c c_c, \\[8pt]
    \dfrac{\partial s_a}{\partial c_a} - \dfrac{\partial s_a}{\partial c_b}, & \text{otherwise (limit branch)},
  \end{cases}
\end{equation}
followed by $\ud\bS = \bN\ud\hat{\bS}\bN^T$ and
$\ud\bP = \ud\bF\,\bS_{\mathrm{iso}} + \bF\ud\bS$.  The partial
derivatives $\partial s_a/\partial c_b$ follow from
$\partial\beta_a/\partial c_b = \beta'_a(\delta_{ab}/(2c_a) - 1/(6c_b))$
with $\beta'_a = \sum_r\mu_r\alpha_r\bar\lambda_a^{\alpha_r}$.  Ogden with
one term $(\mu, 2)$ is the isochoric neo-Hookean model, with two terms
$(2c_1, 2), (-2c_2, -2)$ the Mooney--Rivlin model.

\subsection{Category 3: slightly compressible}
\label{sec:nearly}

The same six materials with $\kappa \gg \mu$ describe rubber-like solids
whose bulk modulus is finite but three or four orders of magnitude above
the shear modulus (\code{nu} close to $\tfrac12$, or \code{kappa} given
directly).  The constitutive law is unchanged, but its numerical treatment
is not.  In the displacement formulation the volumetric term is a stiff
penalty: the discrete field must satisfy $J \approx 1$ nearly pointwise,
which the low-order continuous displacement space cannot do without
suppressing admissible deformations (volumetric locking), and the
pressure recovered from $U'(J)$ oscillates because small errors in
$J$ are multiplied by $\kappa$.  Newton also sees a badly scaled tangent
dominated by the penalty.

The remedy is to treat the pressure as an independent field: replace
$U'(J)$ in \eqref{eq:Pdisp} by an unknown $p$ and impose the volumetric law
weakly,
\begin{equation}
  \bP = \bP_{\mathrm{iso}}(\bF) + p\,J\bF^{-T},
  \qquad
  \int_{\Omega_R} q\,\Bigl(u'(J) - \frac{p}{\kappa}\Bigr)\ud V = 0 \quad \forall q,
  \label{eq:perturbed}
\end{equation}
with $p$ taken from a lower-order space (Section~\ref{sec:mixed}), which the
Taylor--Hood pair satisfies without locking.  The constraint is written
with the normalised pressure $u' = U'/\kappa$ so that it stays finite at
$\kappa = \infty$; for the quadratic law it is the familiar
$J - 1 - p/\kappa$.  For that law the pair is the stationarity of the
perturbed Lagrangian
$\int\Psi_{\mathrm{iso}} + p(J-1) - p^2/(2\kappa)\ud V$, the Legendre
transform of $U$, and the Newton system is symmetric; for the other laws
the constraint is not the derivative of a potential in $(\bu, p)$ jointly,
and the $(2,1)$ block carries the factor $u''(J)$ (Section~\ref{sec:mixed}),
which the FGMRES saddle-point solver accepts.  Eliminating $p$ pointwise
from \eqref{eq:perturbed} recovers \eqref{eq:Pdisp} exactly for every law,
so the mixed and displacement formulations solve the same continuous
problem at any finite $\kappa$; their discrete solutions converge to each
other under refinement, which the tests check at $\nu = 0.45$.  The
pressure unknown is the mean stress $p = \tfrac13\tr\bsig$, tensile
positive, and the $(2,2)$ block $-\bM_p/\kappa$ of the Newton system keeps
it uniquely determined even when the displacement is prescribed on the
whole boundary.  The bulk modulus is resolved from $\nu$ as
$\kappa = 2\mu(1+\nu)/(3(1-2\nu))$ (Section~\ref{sec:moduli}); at
$\nu = 0.4999$ this gives $\kappa/\mu \approx 5000$, the value of the
nearly incompressible Cook input.  The saddle-point solver of
Section~\ref{sec:saddlesolver} is built for this regime: its Schur
complement approximation $\bM_p(1/\mu + 1/\kappa)$ reduces to the
incompressible one as $\kappa \to \infty$.

\subsection{Category 4: incompressible}
\label{sec:incompressible}

In the limit $\kappa \to \infty$ the volumetric energy becomes the
internal constraint $J = 1$ and the pressure a Lagrange multiplier that
does no work:
\begin{equation}
  \bP = \bP_{\mathrm{iso}}(\bF) + p\,J\bF^{-T}, \qquad
  \int_{\Omega_R} q\,(J - 1)\ud V = 0 \quad \forall q .
  \label{eq:incompressible}
\end{equation}
The input selects it with \code{incompressible: true} or \code{nu: 0.5};
the materials report $\kappa = \infty$, the $(2,2)$ block of the Newton
system vanishes, and the mixed formulation is required (the displacement
formulation refuses the input).  With the displacement prescribed on the
whole boundary the multiplier is determined only up to a constant
(Section~\ref{sec:mixed}), so incompressible problems need a traction
boundary somewhere.  The one exception to the mixed requirement is plane
stress (Section~\ref{sec:planestress}): for a thin sheet the constraint
fixes the thickness stretch $\lambda_3 = 1/\det\bF_{2D}$ and the
traction-free face fixes $p = -\sigma_{\mathrm{iso},33}$ pointwise, so the
multiplier is eliminated at the material point and the sheet is a pure
displacement problem.

On an incompressible solution the barred and unbarred invariants coincide
and the six isochoric energies reduce to the classical incompressible
models $\Psi(\Ione, \Itwo)$ or $\Psi(\lambda_1, \lambda_2, \lambda_3)$ with
$\lambda_1\lambda_2\lambda_3 = 1$.  With $\Psi_1 = \partial\Psi/\partial\Ione$
and $\Psi_2 = \partial\Psi/\partial\Itwo$ the Cauchy stress is
\begin{equation}
  \bsig = 2\Psi_1\dev\bB - 2\Psi_2\dev\bB^{-1} + p\,\bI,
  \qquad
  \mu = 2(\Psi_1 + \Psi_2)\big|_{\Ione = \Itwo = 3},
\end{equation}
and the principal stress differences
$\sigma_a - \sigma_b = \lambda_a\partial\Psi/\partial\lambda_a - \lambda_b\partial\Psi/\partial\lambda_b$
are constitutive while the mean stress is fixed by equilibrium and the
boundary conditions.  Much of the literature writes
$\bsig = -\tilde p\,\bI + 2\Psi_1\bB - 2\Psi_2\bB^{-1}$ and calls $\tilde p$
the hydrostatic pressure; its value depends on which isotropic terms an
author absorbs into it, whereas $p = \tfrac13\tr\bsig$ does not, and that is
the unknown used here.  The closed-form uniaxial, equibiaxial, pure-shear
and simple-shear states of these models, for a general $\Psi(\Ione, \Itwo)$
and specialised to the six energies, serve as reference solutions and are
collected in the appendix of the verification manual, which also lists the
inputs that reproduce them.

\subsection{Resolution of the moduli and material regions}
\label{sec:moduli}

The input names a model and its parameters; the factory
(\file{materials.cpp}) resolves the small-strain shear and bulk moduli
$(\mu, \kappa)$ as follows.  The coupled models take $(E, \nu)$ with
$\nu < \tfrac12$.  For a decoupled model $\mu$ is the model's own
small-strain value from the tables above (or $E/(2(1+\nu))$ for the
isochoric neo-Hookean model given $E, \nu$), and the bulk modulus comes
from exactly one of
\begin{equation}
  \kappa \ \text{given}, \qquad
  \kappa = \frac{2\mu(1+\nu)}{3(1 - 2\nu)}\ \text{from } \nu < \tfrac12, \qquad
  \kappa = \infty \ \text{from } \nu = \tfrac12 \text{ or \code{incompressible: true}},
\end{equation}
which places the material in category 2, 3 or 4.  The derived Lam\'e
constant $\lambda = \kappa - \tfrac23\mu$ is used only in the small-strain
checks of the tests.

Materials may vary by element attribute (\code{material.regions}): every
region uses the base model with other parameter values, inheriting the
keys it does not give (a region that gives any bulk key replaces the
base's bulk specification).  The integrators hold a table indexed by
attribute and pick the entry per element; all regions must be
incompressible or none, and no attribute may be covered twice.

%============================================================================
\section{Two-dimensional problems}
\label{sec:planestress}
%============================================================================

\subsection{Plane strain}

By default a two-dimensional mesh is treated as plane strain: the
displacement has two components, $F_{33} = 1$ in \eqref{eq:embed}, and the
three-dimensional material returns the full $\bP$, of which the kernels
contract only the in-plane block with $\Grad\bw$.  The out-of-plane stress
$\sigma_{33}$ is a reaction; it is generally non-zero and is included in
the pressure, the von Mises stress and the stress outputs.  An
incompressible plane-strain state has in-plane stretches
$(\lambda, \lambda^{-1})$.

\subsection{Plane stress}

For a thin sheet loaded in its plane the appropriate reduction is
$\sigma_{i3} = 0$.  The thickness stretch $\lambda_3 = F_{33}$ is then not
prescribed but determined by the material.  \code{plane: stress}
(two-dimensional, displacement formulation) wraps the material in the
adapter of \file{plane\_stress.hpp}: the kernels still pad $\bF$ with
$F_{33} = 1$, the adapter replaces $F_{33}$ by $\lambda_3$ (and drops the
out-of-plane shears) such that $P_{33} = 0$, and returns $\bP$ at the
completed $\bF_c$; its $(3,3)$ entry is zero and the in-plane block is the
plane-stress response.  Two branches:

\paragraph{Incompressible base ($\kappa = \infty$).}
The constraint fixes the thickness pointwise and $\sigma_{33} = 0$ fixes
the mean stress:
\begin{equation}
  \lambda_3 = \frac{1}{\det\bF_{2D}}, \qquad
  p = -\sigma_{\mathrm{iso},33} = -\lambda_3\,P_{\mathrm{iso},33}, \qquad
  \bP = \bP_{\mathrm{iso}}(\bF_c) + p\,\bF_c^{-T}
  \quad (J = 1).
  \label{eq:psinc}
\end{equation}
The Lagrange multiplier is eliminated at the material point: the sheet
problem is a pure displacement problem in the two in-plane components with
no pressure unknown, no inf-sup condition and no locking.  This is the one
way to run an incompressible model in the displacement formulation.

\paragraph{Compressible base.}
$\lambda_3$ solves the scalar equation $r(\lambda_3) = P_{33}(\bF_{2D}, \lambda_3) = 0$
by Newton's method from $\lambda_3 = 1$, with $\ud r/\ud\lambda_3$ from a
dual evaluation of $P_{33}$, steps clipped to $[-\lambda_3/2, \lambda_3]$
to keep the stretch positive, a relative tolerance of $10^{-14}$ and at
most 100 iterations; a non-finite state (an inverted in-plane block during
a line search) returns NaN, which the outer Newton rejects.  When the
adapter is evaluated with dual numbers the converged root is refined by
\emph{one Newton step in dual arithmetic},
\begin{equation}
  \lambda_3^{\mathrm{dual}} = \lambda_3 - \frac{r(\lambda_3)^{\mathrm{dual}} - r(\lambda_3)}{\ud r/\ud\lambda_3},
\end{equation}
which leaves the value unchanged and gives the derivative part
$-(\partial P_{33}/\partial\bF:\ud\bF)/(\partial P_{33}/\partial F_{33})$,
the implicit-function derivative $\ud\lambda_3/\ud\bF$.  The in-plane
tangent obtained by dual seeding is therefore the consistent plane-stress
tangent.  The energy is the base energy at the completed $\bF_c$; since
$P_{33} = 0$ there, its in-plane derivative is the in-plane $\bP$
(envelope theorem).  In the small-strain limit the adapter reproduces the
plane-stress moduli $E/(1-\nu^2)$, $E\nu/(1-\nu^2)$ and the thickness
strain $-\nu/(1-\nu)\,\tr\beps$.

Every quadrature quantity uses the completed $\bF_c$ (so \code{jacobian}
is the true $J$, one for an incompressible sheet), and
\code{thickness\_stretch} $= \lambda_3$ becomes available.  Two cautions:
the idealisation assumes stress uniform through the thickness and no
bending; and because each load step starts with the boundary displaced and
the interior lagging, shear-dominated sheet problems need increments that
shear a boundary layer of elements by no more than a few percent, since the
energy grows like $(\det\bF_{2D})^{-2}$.

\subsection{Axisymmetry}
\label{sec:axisymmetric}

\code{plane: axisymmetric} (either formulation) reads the two-dimensional
mesh as the meridian section of a solid of revolution, $x = r$ and $y = z$,
without torsion: the displacement has the components $(u_r, u_z)$, the
reference position of a material point is $(r, z)$ and its current one
$(r + u_r, z + u_z)$ in the same meridian plane.  The deformation gradient
in the cylindrical basis $(\be_r, \be_z, \be_\phi)$ is the in-plane block
$\bI + \Grad\bu$ of \eqref{eq:embed} completed by the hoop stretch
\begin{equation}
  F_{33} = \frac{r + u_r}{r} = 1 + \frac{u_r}{r},
  \qquad F_{33} \to 1 + \partial_r u_r = F_{rr} \ \text{on the axis},
  \label{eq:hoop}
\end{equation}
the ratio of the current to the reference circumference, with no
out-of-plane shears; on the axis, where $u_r/r$ is indeterminate, the limit
$F_{rr}$ is used (a point is on the axis when $r$ is not positive).  The
volume element of the solid of revolution is $2\pi r\,\ud r\,\ud z$ and a
surface element $2\pi r\,\ud s$, so every integral of the weak forms of
Sections~\ref{sec:momentum} and \ref{sec:mixed}, of the loads of
Section~\ref{sec:loads} and of the mass matrix of Section~\ref{sec:dynamics}
carries the weight $2\pi r$: energies, reactions, resultants and masses are
the totals over the full revolution (per unit radian times $2\pi$), and a
traction or pressure of the input acts on the real surface.  The variation
of $F_{33}$ with respect to a radial nodal displacement is $N_a/r$
($\partial_r N_a$ on the axis), so the B-operator of a displacement dof
carries that hoop term in the $(3, 3)$ slot next to the in-plane entries,
the residual contracts the full $3\times3$ stress with it,
\begin{equation}
  R_{ai} = \int_{\Omega_R} \Bigl[P_{ij}\,\partial_j N_a + \delta_{ir}\,P_{33}\,\frac{N_a}{r}\Bigr] 2\pi r\,\ud A,
\end{equation}
and the material tangent of Section~\ref{sec:ad} is seeded on the four
in-plane entries and the hoop one, five directions instead of four.  The
face kernels (follower pressure, contact, the heat flux of
Section~\ref{sec:thermoelasticity}) complete $\bF$ the same way at their
quadrature points and weight by $2\pi r$; the dead loads and the body force
are multiplied by $2\pi r$ through their coefficients; the reference
density that enters the mass matrix is multiplied likewise; the quadrature
outputs use the completed $\bF$.  Plane strain and three dimensions are
untouched: the axisymmetric branch is a flag of the kernels tested once per
element.  The axis is a boundary of the mesh like any other and $u_r = 0$
is to be prescribed there.  The inputs
\file{rivlin\_cylinder\_axisymmetric.yaml} and
\file{green\_zerna\_sphere\_axisymmetric.yaml} run the inflated cylinder
and sphere of the verification set in this description
(\file{tests/test\_axisymmetric.cpp}).

%============================================================================
\section{Mixed displacement--pressure formulation}
\label{sec:mixed}
%============================================================================

For near- and fully incompressible materials the pressure is introduced as
an independent field.  With $U(J) = \kappa\,u(J)$ the material's volumetric
law (Section~\ref{sec:split}), the two-field residuals are, for all
$\bw \in \mathcal V$ and all $q$ in the pressure space,
\begin{equation}
  \boxed{\;
  \begin{aligned}
  R_u(\bu, p; \bw) &= \int_{\Omega_R}\bigl[\bP_{\mathrm{iso}}(\bF) + p\,J\bF^{-T}\bigr]:\Grad\bw\ud V \\
    &\quad - \int_{\Omega_R}\rho_R\bb\cdot\bw\ud V - \int_{\Gamma_R^N}\bar\bT\cdot\bw\ud A
    + R_{\text{fol}}(\bu; \bw) = 0, \\
  R_p(\bu, p; q) &= \int_{\Omega_R} q\,\Bigl(u'(J) - \frac{p}{\kappa}\Bigr)\ud V = 0 .
  \end{aligned}\;}
  \label{eq:mixedweak}
\end{equation}
For the default quadratic law $u'(J) = J - 1$ and the second equation is
the familiar $\int q\,(J - 1 - p/\kappa)\ud V = 0$; for the quadratic law
alone the pair is the stationarity of the perturbed Lagrangian
\begin{equation}
  \Pi(\bu, p) = \int_{\Omega_R}\Bigl[\Psi_{\mathrm{iso}}(\bF) + p\,(J - 1) - \frac{p^2}{2\kappa}\Bigr]\ud V
  - \Pi_{\mathrm{ext}}(\bu).
  \label{eq:lagrangian}
\end{equation}
For finite $\kappa$ the second equation is the weak form of the volumetric
law $p = U'(J)$, and the mixed and displacement formulations solve the same
continuous problem; for $\kappa = \infty$ it reads $\int q\,u'(J)\ud V = 0$,
the incompressibility constraint $J = 1$ in weak form for every law, and
$p$ is a Lagrange multiplier.  With $\bP_{\mathrm{iso}}:\bF = 0$ the Cauchy stress is
$\bsig = J^{-1}\bP_{\mathrm{iso}}\bF^T + p\bI$, so the unknown $p$ is the
mean stress $\tfrac13\tr\bsig$, tensile positive, and the constitutive
models need only provide the split of Section~\ref{sec:split}.

\paragraph{Spaces.}
Displacement and pressure are discretised on a Taylor--Hood pair,
$[H^1]^d$ of order $k$ for $\bu$ and $H^1$ of order $k - 1$ for $p$, so the
input needs $k \ge 2$.  The affine displacement and the constant pressure
lie in these spaces, which is why homogeneous states are reproduced to
solver tolerance.  Dirichlet conditions act on the displacement only; the
pressure carries no essential conditions.  With the displacement prescribed
on the whole boundary an incompressible pressure is determined only up to a
constant, because
$c\int J\bF^{-T}:\Grad\bw\ud V = c\int_{\partial\Omega_R}\bw\cdot J\bF^{-T}\bN\ud A = 0$
for every admissible $\bw$ (Piola identity); a traction boundary somewhere
fixes it.

\paragraph{Tangent.}
The Newton system is the block system
\begin{equation}
  \begin{bmatrix} \bK & \bBm \\ \bBm_{pu} & -\bM_p/\kappa \end{bmatrix}
  \begin{bmatrix} \Delta\bu \\ \Delta p \end{bmatrix}
  = -\begin{bmatrix} R_u \\ R_p \end{bmatrix},
  \label{eq:saddle}
\end{equation}
with, at each quadrature point,
\begin{gather}
  \bK: \ \Grad\bw:\bA(\bF, p):\Grad\Delta\bu, \qquad
  \bA = \frac{\partial}{\partial\bF}\bigl[\bP_{\mathrm{iso}} + pJ\bF^{-T}\bigr]_{p\ \text{fixed}}, \\
  \bBm: \ q\,(J\bF^{-T}):\Grad\bw, \qquad
  \bBm_{pu}: \ q\,u''(J)\,(J\bF^{-T}):\Grad\Delta\bu, \qquad
  \bM_p: \ q\,\Delta p .
\end{gather}
The $(2,2)$ block vanishes when $\kappa = \infty$.  $\bA$ at fixed $p$ is
again produced by dual seeding of the combined stress; $\bBm$ and
$\bBm_{pu}$ use $\partial J/\partial\bF = J\bF^{-T}$.  For the quadratic
law $u'' = 1$ and $\bBm_{pu} = \bBm^T$, so the system is symmetric; for the
other laws it is not, and the linear solver of
Section~\ref{sec:saddlesolver} is written for the general case.

%============================================================================
\section{Small-strain linear elasticity}
\label{sec:smallstrain}
%============================================================================

\subsection{The same weak form}

Linearising the kinematics about $\bF = \bI$ replaces every strain measure
by the infinitesimal strain and removes the distinction between the
reference and the current configuration: $\Omega_R$ is the domain of
integration, and $\bP$, $\bS$ and $\bsig$ coincide.  With
$\bH = \Grad\bu = \bF - \bI$,
\begin{equation}
  \beps = \sym\bH, \qquad
  \bsig = \mathbb{C}:\beps = 2\mu\,\dev\beps + \kappa\,(\tr\beps)\,\bI
        = \lambda\,(\tr\beps)\,\bI + 2\mu\,\beps, \qquad
  \lambda = \kappa - \tfrac23\mu ,
  \label{eq:hooke}
\end{equation}
with the energy $\Psi = \mu\,\dev\beps:\dev\beps + \tfrac12\kappa(\tr\beps)^2$.
Because $\bsig$ is symmetric, $\bsig:\beps(\bw) = \bsig:\Grad\bw$, and the
weak form of linear elasticity,
\begin{equation}
  \int_{\Omega_R}\bsig(\beps(\bu)):\beps(\bw)\ud V
  = \int_{\Omega_R}\bsig:\Grad\bw\ud V
  = \int_{\Omega_R}\rho_R\bb\cdot\bw\ud V + \int_{\Gamma_R^N}\bar\bT\cdot\bw\ud A ,
  \label{eq:linearweak}
\end{equation}
is the total Lagrangian form \eqref{eq:weak} with the flux
$\bP(\bF) := \bsig(\sym(\bF - \bI))$.  It is therefore not a new weak form
and not a new integrator: \code{linear\_elastic} is a material whose
\code{PK1(F)} returns \eqref{eq:hooke}.  The element residual is
$\int\sigma_{ij}N_{a,j}$, the tangent
$\partial\bP/\partial\bF = \mathbb{C}$,
$C_{ijkl} = \lambda\delta_{ij}\delta_{kl} + \mu(\delta_{ik}\delta_{jl} + \delta_{il}\delta_{jk})$,
is constant and has the minor symmetries, and the element matrix
$\int N_{a,j}C_{ijkl}N_{b,l}$ is the classical $\bm B^T\mathbb{C}\bm B$
stiffness (it equals MFEM's \code{ElasticityIntegrator} to round-off; the
tangents of the compressible hyperelastic models at $\bu = \bm 0$ are this
matrix).  The model is not objective under finite rotations, by
construction; it is invariant under infinitesimal ones, $\bP(\bI + \bm W) = \bm 0$
for skew $\bm W$.  The input keys are those of \code{iso\_neo\_hookean}:
\code{mu}, or \code{E} and \code{nu}, and the bulk modulus from exactly one
of \code{kappa}, \code{nu}, \code{incompressible}; there is no volumetric
law to choose.

\subsection{What depends on the kinematics}

The flux contract of the kernels is the same for both kinematics.  What
differs lies outside the flux, and is switched by a compile-time trait of
the material (\file{materials/kinematics.hpp}; the plane-stress adapter
inherits the trait of its base):
\begin{center}
\small
\renewcommand{\arraystretch}{1.15}
\begin{tabular}{@{}>{\raggedright\arraybackslash}p{4.4cm} >{\raggedright\arraybackslash}p{4.9cm} >{\raggedright\arraybackslash}p{5.6cm}@{}}
\toprule
 & Finite strain & Small strain \\
\midrule
Cauchy stress (outputs, von Mises) & $\bsig = J^{-1}\bP\bF^T$ & $\bsig = \bP$ \\
Volume ratio (\code{jacobian}) & $J = \det\bF$ & $1 + \tr\beps$ \\
\code{strain} & Green--Lagrange $\bE = \tfrac12(\bF^T\bF - \bI)$ & $\beps = \sym(\bF - \bI)$ \\
Volume measure of the mixed form and its gradient & $J$, \ $\partial J/\partial\bF = J\bF^{-T}$ & $\theta = 1 + \tr\beps$, \ $\partial\theta/\partial\bF = \bI$ \\
Incompressible plane stress & $\lambda_3 = 1/\det\bF_{2\mathrm D}$, $\bP = \bP_{\mathrm{iso}} + p\,\bF^{-T}$ & $\varepsilon_{33} = -(\varepsilon_{11} + \varepsilon_{22})$, $\bP = \bP_{\mathrm{iso}} + p\,\bI$ \\
Follower pressure & a load of its own (Section~\ref{sec:follower}) & identical to the dead pressure; rejected as an input error \\
Moment arms of the reactions & current positions $\bX + \bu$ & reference positions $\bX$ \\
\bottomrule
\end{tabular}
\end{center}
A stress formed with the finite-strain push-forward would be wrong by
$O(|\Grad\bu|)$ and not symmetric; with reference moment arms the moment
balance of the reactions closes exactly, as the equilibrium equations of the
linear theory are written on the reference configuration.  The outputs
\code{pk1\_stress} and \code{cauchy\_stress} coincide, the
\code{deformation\_gradient} is $\bI + \Grad\bu$, and the
\code{thickness\_stretch} is $1 + \varepsilon_{33}$.

\subsection{Two-dimensional problems}

Plane strain is the padded $F_{33} = 1$, i.e. $\varepsilon_{33} = 0$, with
$\sigma_{33} = \lambda(\varepsilon_{11} + \varepsilon_{22})$ carried in the
outputs.  Under plane stress the compressible branch of the adapter of
Section~\ref{sec:planestress} applies unchanged: its scalar Newton iteration
on $P_{33} = 0$ is exact after one step and returns
$\varepsilon_{33} = -\tfrac{\nu}{1-\nu}(\varepsilon_{11} + \varepsilon_{22})$,
and the dual refinement returns the plane-stress moduli,
$\sigma_{\alpha\beta} = \lambda^*\varepsilon_{\gamma\gamma}\delta_{\alpha\beta} + 2\mu\,\varepsilon_{\alpha\beta}$
with $\lambda^* = 2\lambda\mu/(\lambda + 2\mu)$, that is $E/(1-\nu^2)$ and
$E\nu/(1-\nu^2)$.  For $\nu = \tfrac12$ the thickness strain and the mean
stress are eliminated in closed form, $\lambda^* = 2\mu$, and no pressure
unknown is needed.

\subsection{Mixed formulation (Herrmann)}

With the volume measure $\theta = 1 + \tr\beps$ and
$\partial\theta/\partial\bF = \bI$ in place of $J$ and $J\bF^{-T}$, the
residuals \eqref{eq:mixedweak} become
\begin{equation}
  R_u = \int_{\Omega_R}\bigl[2\mu\,\dev\beps + p\,\bI\bigr]:\Grad\bw\ud V - \text{loads}, \qquad
  R_p = \int_{\Omega_R} q\,\Bigl(\tr\beps - \frac{p}{\kappa}\Bigr)\ud V ,
  \label{eq:herrmann}
\end{equation}
Herrmann's formulation of linear elasticity \cite{herrmann}: the material provides
$\bP_{\mathrm{iso}} = 2\mu\,\dev\beps$ and the volumetric law
$u(\theta) = \tfrac12(\theta - 1)^2$, and $\kappa = \infty$ gives the
incompressible problem $\tr\beps = 0$ in weak form.  The blocks of
\eqref{eq:saddle} are $\bBm: q\,\dive\bw$, $\bBm_{pu} = \bBm^T$, and $\bK$
carries no pressure term, so the system is symmetric, independent of the
state, and $\bK$ is positive semi-definite; the saddle-point solver of
Section~\ref{sec:saddlesolver} applies unchanged.  As for the finite-strain
materials the constraint holds weakly: the pointwise $\tr\beps$ of an
incompressible solution is $O(h^2)$, not zero.

\subsection{Solution of a linear problem}
\label{sec:linearsolve}

The residual is affine in the unknown, so the first Newton step is the exact
solve: one assembly, one preconditioner setup, one Krylov solve.  Four
points follow.
\begin{itemize}[nosep]
  \item \emph{The residual has a round-off floor that can lie above the
    Newton tolerance.}  The residual $\bK\bu - \bm f$ is evaluated with an
    error of the order of $\epsilon_{\mathrm{mach}}\|\bK\|\|\bu\|$, that is
    $\epsilon_{\mathrm{mach}}\|\bK\|\|\bu\|/\|\bm f\|$ relative to
    $\|R_0\|$.  Bending makes this ratio large: measured values are
    $4\times10^{-11}$ for Cook's membrane, $8\times10^{-11}$ for a 10:1
    cantilever and $1.6\times10^{-9}$ for a 20:1 beam, and tightening the
    Krylov tolerance does not move them.  A nonlinear problem passes such a
    tolerance by quadratic convergence; a linear one lands on the floor and
    stays there, and a second step then fails its line search.  A problem
    therefore declares itself linear (\code{QuasiStaticProblem::IsLinear}:
    a small-strain material with dead loads), and for such a problem a second
    or later Newton step that fails its line search, or reduces the residual
    by less than a factor of ten, is taken as having met the floor: the
    state is accepted (\code{NewtonReport::at\_floor}), provided every
    linear solve converged.  Steps that still reduce the residual tenfold
    (an inexact linear solver, i.e. iterative refinement) continue as usual,
    and nothing changes for a nonlinear problem.
  \item \emph{Tolerances.}  \code{newton.rtol} only accepts the linear solve;
    the accuracy of the solution is set by \code{linear.rtol}.  The inputs
    use $10^{-8}$ and $10^{-13}$.  The mixed formulation typically takes a
    second step, which refines the first.
  \item \emph{The tangent predictor is excluded.}  After an exact predictor
    Newton would start at the floor, where $\text{rtol}\,\|R_0\|$ is
    meaningless; \code{solver.predictor: tangent} is an input error for a
    small-strain material.
  \item \emph{The Jacobian is assembled once for a load path.}  It depends
    neither on the state nor on the pseudo-time (the essential degrees of
    freedom are fixed when the loads are finalised), so a problem that
    declares itself linear returns the same assembled matrix from every
    Jacobian request until its boundary conditions change.  The physics
    shares an \emph{operator stamp} with the linear solver it creates, a
    counter incremented at every assembly; a solver that is handed the same
    operator with an unchanged stamp keeps its setup, the AMG hierarchy and,
    in the mixed formulation, the augmented blocks of
    Section~\ref{sec:saddlesolver}.  The address of the matrix alone would
    not identify it, since a reassembled matrix can occupy the address of the
    old one.  A nonlinear problem stamps at every request, so nothing changes
    for it.  The results are identical bit for bit to reassembling; over ten
    steps the time spent in assembly and setup drops by a factor of seven to
    nine.  The Krylov solves remain and dominate these runs, so the overall
    gain is about $1.5$ for the displacement formulation and negligible for
    the mixed one.
\end{itemize}
Forming $\bF = \bI + \bH$ and subtracting $\bI$ again leaves an absolute
error of $\epsilon_{\mathrm{mach}}$ in $\beps$, a relative floor of
$\epsilon_{\mathrm{mach}}/|\Grad\bu|$ in the stress.  Since the solution
scales with the load, inputs choose amplitudes with
$|\Grad\bu| \gtrsim 10^{-4}$ and scale the results, never the reverse; the
``linear limit'' runs of the hyperelastic models, which must keep
$|\Grad\bu|$ small, do not have that freedom.

%============================================================================
\section{Elastodynamics}
\label{sec:dynamics}
%============================================================================

A \code{dynamics} block in the input adds the inertial term to either weak
form.  Without it nothing of this section exists: the quasi-static analysis
is unchanged to the last digit.

\subsection{Strong and weak form; the constant mass matrix}

With $t$ now the physical time, balance of linear momentum on the reference
configuration and its weak form are
\begin{equation}
  \rho_R\ddot\bu = \Dive\bP + \rho_R\bb \ \text{in } \Omega_R, \qquad
  \int_{\Omega_R}\rho_R\ddot\bu\cdot\bw\ud V + R(\bu; \bw) = 0
  \quad \forall \bw \in \mathcal V ,
  \label{eq:dynweak}
\end{equation}
with $R$ the quasi-static residual \eqref{eq:weak} (internal minus
external) and initial data $\bu(\cdot, 0) = \bu_0$,
$\dot\bu(\cdot, 0) = \bm v_0$.  In total Lagrangian form the inertial term is
integrated over the \emph{fixed} mesh with the \emph{reference} density: mass
conservation, $\rho_R = J\rho$, is built in, and the mass matrix
\begin{equation}
  M_{ai,bj} = \delta_{ij}\int_{\Omega_R}\rho_R N_a N_b\ud V
\end{equation}
is constant.  It is assembled once, the same for finite and small strain and
for every material; it is the consistent mass, with $\rho_R$ piecewise
constant by element attribute (\code{material.rho0} and the regions'), the
same coefficient that multiplies the body force.  In two dimensions it is
the mass per unit reference thickness, also under plane stress, where the
current thickness changes and the mass does not.  The semi-discrete system
is
\begin{equation}
  \bM\ddot{\bm u} + \bm S(\bm u, t) = \bm 0 ,
  \label{eq:semidiscrete}
\end{equation}
where $\bm S$ is the static residual vector of the wrapped problem (its
\code{Mult}) and $\bK = \partial\bm S/\partial\bm u$ its Jacobian.

\subsection{Time integration in displacement form}

One family of implicit one-step schemes is implemented
(\file{solvers/time\_integration.hpp}), with the displacement of the new
time level as the unknown, so that prescribed displacements, Newton, the
line search and the linear solvers are those of the quasi-static analysis.
With $(\bm u_n, \bm v_n, \bm a_n)$ known and $\Delta t = t_{n+1} - t_n$,
\begin{equation}
\begin{aligned}
  \bm u^* &= \bm u_n + \Delta t\,\bm v_n + \Delta t^2(\tfrac12 - \beta)\bm a_n, &
  \bm a(\bm u) &= \frac{\bm u - \bm u^*}{\beta\Delta t^2}, \\
  \bm v^* &= \bm v_n + \Delta t(1 - \gamma)\bm a_n, &
  \bm v(\bm u) &= \bm v^* + \gamma\Delta t\,\bm a(\bm u),
\end{aligned}
\end{equation}
the balance is imposed between the time levels with \emph{interpolated
forces} (the form of the HHT operator of general-purpose codes),
\begin{equation}
  \bM\bigl[(1 - \alpha_m)\bm a(\bm u) + \alpha_m\bm a_n\bigr]
  + (1 - \alpha_f)\,\bm S(\bm u, t_{n+1}) + \alpha_f\,\bm S_n = \bm 0,
  \qquad \bm S_n = \bm S(\bm u_n, t_n) .
  \label{eq:genalpha}
\end{equation}
Interpolating the forces rather than the state means $\bm S$ is only ever
evaluated at $(\bm u_{n+1}, t_{n+1})$: the load set, the Dirichlet data and
the follower loads keep one time.  The parameter sets are
\begin{center}
\small
\begin{tabular}{@{}l l l l l@{}}
\toprule
\code{scheme} & $\alpha_m$ & $\alpha_f$ & $\beta$ & $\gamma$ \\
\midrule
\code{newmark} & 0 & 0 & given ($\tfrac14$) & given ($\tfrac12$) \\
\code{hht} & 0 & $\alpha \in [0, \tfrac13]$ & $\tfrac14(1 + \alpha)^2$ & $\tfrac12 + \alpha$ \\
\code{generalized\_alpha} & $\dfrac{2\rho_\infty - 1}{\rho_\infty + 1}$ & $\dfrac{\rho_\infty}{\rho_\infty + 1}$ & $\tfrac14(1 - \alpha_m + \alpha_f)^2$ & $\tfrac12 - \alpha_m + \alpha_f$ \\
\bottomrule
\end{tabular}
\end{center}
(Newmark~\cite{newmark}; Hilber, Hughes and Taylor~\cite{hht}; Chung and
Hulbert~\cite{chunghulbert}).  All are second-order accurate with
$\gamma = \tfrac12 - \alpha_m + \alpha_f$ and, for a linear problem,
unconditionally stable; \code{newmark} with $\gamma > \tfrac12$ is first
order, and with $2\beta < \gamma$ only conditionally stable (the run header
says so).  $\rho_\infty \in [0, 1]$ is the spectral radius of the
amplification matrix at infinite frequency: $1$ means no numerical
dissipation, $0$ annihilates in a few steps whatever the step does not
resolve.  The default \code{newmark} pair $(\tfrac14, \tfrac12)$ is the
trapezoidal rule, which for a linear problem conserves the energy exactly
and lengthens a period by the factor $(\omega\Delta t/2)/\arctan(\omega\Delta
t/2) \approx 1 + (\omega\Delta t)^2/12$.  For a nonlinear problem it is not
unconditionally stable, and what a mesh cannot resolve behind a wave front
rings for ever; \code{generalized\_alpha} with $\rho_\infty$ between $0.8$
and $0.9$ removes the unresolved part at little cost to the resolved one.
It dissipates in a norm of its own, not step by step in the physical
energy, which may rise slightly between two steps but does not exceed its
initial value.

\paragraph{The step equation.}
Divided by $1 - \alpha_f$, \eqref{eq:genalpha} reads
\begin{equation}
  \boxed{\;
  \bm G(\bm u) = \bm S(\bm u, t_{n+1}) + c_M\bM(\bm u - \bm u^*) + \bm h_n = \bm 0,
  \qquad
  \frac{\partial\bm G}{\partial\bm u} = \bK(\bm u) + c_M\bM ,\;}
  \label{eq:step}
\end{equation}
\begin{equation}
  c_M = \frac{1 - \alpha_m}{(1 - \alpha_f)\,\beta\Delta t^2}, \qquad
  \bm h_n = \frac{\alpha_m\bM\bm a_n + \alpha_f\bm S_n}{1 - \alpha_f} .
\end{equation}
The static operator enters with unit weight: inertia is a linear spring
$c_M\bM$ plus a known history load.  This is why it needs no kernel of its
own.  \code{DynamicSolidProblem} (\file{physics/dynamic\_solid\_problem.hpp})
wraps any solid problem and \emph{is} equation~\eqref{eq:step} behind the
interface the stepper already drives; a linear problem stays linear (the
floor logic of Section~\ref{sec:linearsolve} applies unchanged), and the
consistent Jacobian $\bK + c_M\bM$ keeps Newton quadratic (measured order
$1.9$ at displacements of $0.75$ of the body's size).  After convergence,
$\bm a_{n+1} = \bm a(\bm u_{n+1})$ and $\bm v_{n+1} = \bm v(\bm u_{n+1})$.
$\bM$ is kept twice: whole, for the residual, the energies and the
reactions (it couples the acceleration of a prescribed boundary into the
free rows), and with the essential rows and columns zeroed for the
Jacobian.

\paragraph{Initial state.}
$\bu_0$ and $\bm v_0$ are given as expressions of the reference coordinates
(zero by default); the Dirichlet data of $t = 0$ overwrite $\bu_0$ on the
prescribed degrees of freedom, with a warning if that changes it.  The
initial acceleration is the consistent one,
$\bM\bm a_0 = -\bm S(\bm u_0, 0)$ on the free degrees of freedom and zero
on the prescribed ones, with the loads at $t = 0$ taken as their \emph{right
limit}: a load that is switched on at $t = 0$ and held is on from the first
instant and enters $\bm a_0$.

\paragraph{Prescribed motion and reactions.}
Dirichlet data are imposed on $\bm u_{n+1}$; velocity and acceleration of
those degrees of freedom follow from the update formulas like everywhere
else.  The reaction of a Dirichlet entry is the full balance
$\bm S_{\text{full}}(\bm u_{n+1}, t_{n+1}) + \bM\bm a_{n+1}$ on its rows: a
support force includes the inertia it carries.

\paragraph{Energy balance.}
The kinetic energy is $\tfrac12\bm v\cdot\bM\bm v$, the internal energy that
of Section~\ref{sec:post}, and the external work is accumulated by the
trapezoidal rule over the dead loads and over the supports,
$\tfrac12(\bm f_n + \bm f_{n+1} + \bm r_n + \bm r_{n+1})\cdot(\bm u_{n+1} -
\bm u_n)$ with $\bm r$ the reactions on the prescribed rows.  For a linear
problem and the trapezoidal rule, kinetic plus internal energy minus
external work is constant \emph{exactly} (it is an algebraic identity of the
scheme), which makes the balance a test of the implementation and not only
an output; for a nonlinear problem its error is $O(\Delta t^2)$.  Follower
pressures are not dead loads and are not counted.

\paragraph{A linear problem.}
$\bK + c_M\bM$ is constant while $\Delta t$ is, so it is formed, and the
AMG hierarchy built, once per run (the operator stamp of
Section~\ref{sec:linearsolve}, now of the sum).  Equal steps
$t_{\text{final}}k/n$ differ in the last digits of $t_{k+1} - t_k$; the
increment of the previous step is kept when the new one agrees with it to
$10^{-12}$, which also makes a run independent of how its time grid was
formed.  Such a run is bound by its residual evaluations (three per step
against one Krylov solve); the full static residual that reactions and the
external work need is evaluated only when one of them is asked for.

\subsection{Mixed formulation: a differential-algebraic system}
\label{sec:dynmixed}

In the $u$--$p$ formulation the pressure carries no inertia: $\bM$ acts on
the displacement block, the constraint row is that of the static problem at
$t_{n+1}$, unweighted, and $\bm h_n$ has no pressure block.  The Jacobian is
that of Section~\ref{sec:mixed} with $\bK + c_M\bM$ in place of $\bK$.  Two
properties follow.
\begin{itemize}[nosep]
  \item \emph{The pressure mode.}  The pressure force is part of $\bm S_u$
    and is interpolated with it, so an error $e_n$ in $p_n$ returns as
    $e_{n+1} = -\alpha_f/(1 - \alpha_f)\,e_n = -\rho_\infty e_n$: a
    sign-alternating mode that decays only for $\rho_\infty < 1$ (measured:
    two runs that differ in $p_0$ alone differ at every later step by
    $0.6000$ times the step before, for $\rho_\infty = 0.6$).  The trapezoidal
    rule has the same factor $1$ through the acceleration.  A mixed dynamic
    analysis therefore uses a dissipative scheme, and the run header warns
    otherwise.  The initial pressure is that of $\bu_0$ at finite $\kappa$
    (one solve with the constraint row) and zero for an incompressible
    material, whose true initial pressure would follow from a constraint on
    the acceleration that is not enforced; the start-up error decays with
    the same factor.
  \item \emph{The pressure carries the tolerance of the solves multiplied by
    $c_M$.}  It balances $\bM\bm a$, and $\bm a = (\bm u - \bm u^*)/(\beta
    \Delta t^2)$ amplifies the error of a solve by $O(\Delta t^{-2})$.  With
    \code{newton.rtol} $10^{-10}$ this noise, not the time integration,
    limits the pressure at small steps; tighten the tolerances when the
    pressure matters.
\end{itemize}
With $u$ and $p$ held exactly by the spaces, both are second order in
$\Delta t$, for $\nu = 0.4999$ and for $\kappa = \infty$.  The linear solver
is extended for the mass-dominated displacement block
(Section~\ref{sec:saddlesolver}).

%============================================================================
\section{Finite viscoelasticity}
\label{sec:viscoelasticity}
%============================================================================

Listing Maxwell branches under a decoupled model (\code{material.branches})
makes that model the equilibrium branch of a finite viscoelastic material:
a spring in parallel with $n$ Maxwell elements, each a neo-Hookean spring
of modulus $G_i$ in series with a dashpot of relaxation time $\tau_i$.  It
is the model of the finite viscoelasticity chapter of Anand's coupled
theories and of its FEniCSx codes \cite{anandbook} (the VHB~4910
parameters of their examples are those of Wang et al.\ \cite{wang2016}),
and the first history-dependent material of the library.  This section
gives the model, its time integration, and what a history changes in the
kernels and in the stepping.

\subsection{Constitutive model}

Each branch $i$ carries a multiplicative decomposition of its own,
$\bF = \bF^{e}_i\bF^{v}_i$, with the viscous right Cauchy--Green tensor
$\bC^v_i = \bF^{vT}_i\bF^v_i$ as its internal variable: symmetric, of unit
determinant (the viscous flow is isochoric) and equal to $\bI$ at rest.
With $\bC = \bF^T\bF$ and $\bar\bC = J^{-2/3}\bC$ the free energy per unit
reference volume is
\begin{equation}
  \Psi = \Psi_{\mathrm{iso}}(\bar\bC) + \sum_{i=1}^n \frac{G_i}{2}\bigl[\tr(\bar\bC\,\bC^{v-1}_i) - 3\bigr] + U(J),
  \label{eq:viscoenergy}
\end{equation}
where $\Psi_{\mathrm{iso}}$ and $U$ are those of the equilibrium branch,
any model of Section~\ref{sec:classification} with a split;
$\tr(\bar\bC\,\bC^{v-1}_i)$ is the first invariant of the elastic part
$\bC^e_i = \bF^{v-T}_i\bar\bC\bF^{v-1}_i$ of the branch, so each branch is
a neo-Hookean spring on its own elastic stretch.  The first
Piola--Kirchhoff stress is the derivative at fixed $\bC^v_i$,
\begin{equation}
  \bP = \bP_{\mathrm{iso}}(\bF) + \sum_i G_i J^{-2/3}\Bigl[\bF\bC^{v-1}_i - \tfrac13(\bC:\bC^{v-1}_i)\,\bF^{-T}\Bigr] + U'(J)\,J\bF^{-T},
  \qquad
  \bsig^{\mathrm{neq}}_i = \frac{G_i}{J}\dev\bigl(\bFb\,\bC^{v-1}_i\bFb^T\bigr),
  \label{eq:viscostress}
\end{equation}
deviatoric in the Kirchhoff sense like every isochoric stress of the
library: the branches join $\bP_{\mathrm{iso}}$, while the volumetric law,
the mixed constraint and the pressure field remain those of the equilibrium
branch.  The viscous flow of a branch is driven by its non-equilibrium
Mandel stress $\bM^e_i = G_i\dev\bC^e_i$: the viscous stretching is
$\bm D^v_i = \bM^e_i/(2\eta_i)$ with the viscosity $\eta_i = G_i\tau_i$,
which in terms of the internal variable
($\dot\bC^v_i = 2\bF^{vT}_i\bm D^v_i\bF^v_i$) reads
\begin{equation}
  \dot\bC^v_i = \frac{1}{\tau_i}\Bigl[\bar\bC - \tfrac13\tr(\bar\bC\,\bC^{v-1}_i)\,\bC^v_i\Bigr] .
  \label{eq:viscoflow}
\end{equation}
Its trace term is what keeps $\det\bC^v_i = 1$, since
$\tr(\bC^{v-1}_i\dot\bC^v_i) = 0$.  In the small-strain limit, with
$\bC^v_i = \bI + 2\beps^v_i$ and $\bar\bC = \bI + 2\dev\beps$,
\eqref{eq:viscoflow} is the Maxwell element
$\tau_i\dot\beps^v_i = \dev\beps - \beps^v_i$ and the branch stress is
$2G_i(\dev\beps - \beps^v_i)$: the instantaneous shear modulus is
$\mu + \sum_i G_i$, the relaxed one $\mu$, and a held strain relaxes with
the times $\tau_i$.  The equilibrium branch alone is the hyperelastic
material of the input without the key; every branch adds six internal
variables per quadrature point.

\subsection{Time integration of the internal variables}

Over a step from $t_n$ to $t_{n+1} = t_n + \Delta t$ the reference's codes
update each branch by
\begin{equation}
  \bC^v_{i,n+1} = \frac{\bm A_i}{(\det\bm A_i)^{1/3}}, \qquad
  \bm A_i = \bC^v_{i,n} + \frac{\Delta t}{\tau_i}\,\bar\bC_{n+1},
  \label{eq:viscoupdate}
\end{equation}
the implicit Euler step of \eqref{eq:viscoflow} in which the trace term, a
scalar multiple of $\bC^v_i$, is replaced by the projection onto unit
determinant that it stands for in the continuous equation.  The update is
first-order accurate in $\Delta t$, symmetric and unimodular by
construction, and this code uses it as it stands, so that both codes
integrate the same discrete model.  The stress of the step is
\eqref{eq:viscostress} at $\bF_{n+1}$ with $\bC^v_{i,n+1}$ from
\eqref{eq:viscoupdate}: a material function of $\bF$ with the accepted
history $\bC^v_{i,n}$ as a parameter.  The dual-number tangent of
Section~\ref{sec:ad} differentiates through the update, so the consistent
tangent $\partial\bP/\partial\bF$ at fixed $\bC^v_{i,n}$ includes
$\partial\bC^v_{i,n+1}/\partial\bF$ with no further code (the energy
\eqref{eq:viscoenergy} is not a potential for that stress once
$\bC^v_{i,n+1}$ depends on $\bF$; the tangent is that of the residual, as
in the reference).  With $\Delta t = 0$ the update returns the accepted
value and the stress is that of the accepted state; as
$\Delta t/\tau_i \to \infty$ the branch relaxes to $\bC^v_i = \bar\bC$ and
carries no stress.  In a held small strain the branch stress decays as
$(1 + \Delta t/\tau_i)^{-n}$ over $n$ steps, the implicit Euler image of
$e^{-t/\tau_i}$.

\subsection{Histories in the kernels and in the stepping}

The material (\file{kernels/materials/viscoelastic.hpp}, the class template
\code{Viscoelastic<Base>} over the six decoupled models) evaluates its
stresses and energies from $\bF$, the accepted history block $h$ of the
point (six numbers per branch, the components $xx, yy, zz, xy, yz, xz$ of
$\bC^v_i$) and $\Delta t$, and writes the updated block with \code{Update}.
The kernels keep their stateless contract: a per-point view
(\file{kernels/history\_bound.hpp}) binds the material to the block and to
$\Delta t$ and offers $\bP(\bF)$, $\Psi(\bF)$, $\bP_{\mathrm{iso}}(\bF)$
and $\Psi_{\mathrm{iso}}(\bF)$ to the quadrature-point functions of
Section~\ref{sec:fem}; everything that does not depend on the history
(the volumetric law, $\kappa$, the small-strain shear modulus that scales
the saddle-point preconditioner) is read from the material itself.  The
blocks live in a \code{HistoryField} (\file{kernels/history\_field.hpp}) on
the kernels' quadrature rule, order $2k + 3$ on each element's geometry,
with an accepted and a working copy and the length of the step under way;
a region with fewer branches uses the head of its block.  Both formulations
follow one protocol.  \code{SetLoadFactor(t)} makes $\Delta t = t - t_n$
the length of the step under way, $t_n$ being the time of the last accepted
step.  The accept hook of the stepper (Section~\ref{sec:loadstepping}),
called once per converged increment, evaluates \eqref{eq:viscoupdate} at
every quadrature point from the converged displacement and commits it,
after which $\Delta t = 0$ until the next step begins, so that reactions,
energies and output fields evaluated between steps carry the stress of the
accepted state.  A rejected increment leaves the history untouched, so
bisection needs nothing else, and \code{ResetHistory} restores the rest
state.  In a dynamic analysis the decorator of Section~\ref{sec:dynamics}
advances the wrapped problem's history last in its own accept hook, after
the static residual of the step has served the balance and the external
work: inertia and branches are independent features of one problem.  A
history-dependent material needs a physical time, so the input must carry a
\code{time} block (Section~\ref{sec:loadstepping}) or a \code{dynamics}
block; the pseudo-time of a plain quasi-static analysis has no rate.

%============================================================================
\section{Finite thermoelasticity}
\label{sec:thermoelasticity}
%============================================================================

A \code{thermal} block under a decoupled model makes it a thermoelastic
material and the temperature $\theta$ a third unknown next to the
displacement and the pressure: the model of the finite thermoelasticity
chapter of Anand's coupled theories and of its FEniCSx codes
\cite{anandbook}, whose six examples are reproduced in
\file{apps/input/anand\_coupled\_theories/finite\_thermoelasticity/}.  This
section gives the model, the heat equation and its time integration, the
coupled weak form and its tangents, and where the pieces live.

\subsection{Constitutive model}

The free energy per unit reference volume is
\begin{equation}
  \psi(\bF, \theta) = s(\theta)\,\Psi_{\mathrm{iso}}(\bFb) + \kappa\,u\!\left(\frac{J}{J_\theta}\right) - c_v\Bigl[\theta\ln\frac{\theta}{\theta_0} - (\theta - \theta_0)\Bigr],
  \qquad
  J_\theta = e^{3\alpha(\theta - \theta_0)},
  \label{eq:thermoenergy}
\end{equation}
with $\Psi_{\mathrm{iso}}$ and $u$ those of the base model
(Section~\ref{sec:classification}), $s(\theta) = \theta/\theta_0$ for an
\emph{entropic} elastomer, whose shear modulus is proportional to the
temperature (the default; $s = 1$ turns it off), $J_\theta$ the stress-free
volume ratio of a free thermal expansion with the linear coefficient
$\alpha$, $c_v$ the heat capacity per unit reference volume and $\theta_0$
the reference temperature, which is also the initial one.  The stress at
fixed $\theta$ is
\begin{equation}
  \bP = s(\theta)\,\bP_{\mathrm{iso}}(\bF) + U'(J, \theta)\,J\bF^{-T},
  \qquad
  U'(J, \theta) = \frac{\kappa\,u'(J_m)}{J_\theta}, \quad J_m = \frac{J}{J_\theta},
  \label{eq:thermostress}
\end{equation}
so that a free expansion to $J = J_\theta$ is stress free for every
volumetric law, and for the logarithmic law $U' = \kappa[\ln J - 3\alpha(\theta - \theta_0)]/J$,
the reference's pressure.  In the coupled formulation the volumetric term
is replaced by the pressure field, $\bP = s(\theta)\bP_{\mathrm{iso}} + pJ\bF^{-T}$,
with the constraint
\begin{equation}
  \frac{u'(J_m)}{J_\theta} - \frac{p}{\kappa} = 0
  \qquad (\kappa = \infty:\ u'(J_m) = 0,\ \text{i.e.}\ J = J_\theta),
  \label{eq:thermoconstraint}
\end{equation}
the mixed constraint of Section~\ref{sec:mixed} on the mechanical volume
ratio $J_m$, with the tangents $u''(J_m)/J_\theta^2$ in $J$ and
$-3\alpha[u''(J_m)J_m + u'(J_m)]/J_\theta$ in $\theta$.  The entropy per
unit reference volume is $\eta = -\partial\psi/\partial\theta$; the
thermoelastic coupling of the heat equation is carried by its dependence on
the deformation, through the second Piola--Kirchhoff stress
$\bS = \bF^{-1}\bP$ of the displacement form \eqref{eq:thermostress}:
\begin{equation}
  \bM = \frac{\partial\bS}{\partial\theta} = \bF^{-1}\frac{\partial\bP}{\partial\theta},
  \qquad
  \frac{\partial\bP}{\partial\theta} = s'(\theta)\,\bP_{\mathrm{iso}}(\bF) + \frac{\partial U'}{\partial\theta}\,J\bF^{-T},
  \label{eq:thermoM}
\end{equation}
an explicit function of $(\bF, \theta)$ that the material evaluates with
any scalar type (for the entropic model with the logarithmic law it is the
reference's $J^{-2/3}(G_0/\theta_0)\zeta(\bI - \tfrac13\tr\bC\,\bC^{-1}) - 3\kappa\alpha\,\bC^{-1}$).
Heat is conducted by Fourier's law with the spatial conductivity $k$, whose
referential flux is
\begin{equation}
  \bQ = -k\,J\,\bC^{-1}\Grad\theta .
  \label{eq:fourier}
\end{equation}
The material is the adapter \code{Thermoelastic<Base>}
(\file{kernels/materials/thermoelastic.hpp}) over the six decoupled models,
held in the variant \code{ThermoMaterial}, with the parameters
\code{theta0}, \code{alpha}, \code{c\_v}, \code{k} and \code{entropic}; a
region may change every one of them but $\theta_0$, which is the one
reference temperature of the problem.  Maxwell branches, the small-strain
model and the coupled compressible models cannot carry a temperature.

\subsection{The heat equation and its time integration}

With the entropy as above the balance of energy in the reference
configuration reads
\begin{equation}
  c_v\,\dot\theta = \tfrac12\,\theta\,\bM : \dot\bC - \Div\bQ,
  \label{eq:heat}
\end{equation}
the second term being the thermoelastic (Gough--Joule) heating: an entropic
elastomer warms when stretched and contracts when heated under load.  Over
a step $t_n \to t_{n+1} = t_n + \Delta t$ of the \code{time} block the
equation is integrated by the implicit Euler method, the reference's
scheme,
\begin{equation}
  c_v(\theta - \theta_n) - \tfrac12\,\theta\,\bM(\bF, \theta):(\bC - \bC_n) + \Delta t\,\Div\bQ(\bF, \theta) = 0
  \label{eq:heatstep}
\end{equation}
at the end of the step, with the accepted $\bC_n$ and $\theta_n$ of every
quadrature point held in a \code{HistoryField} of seven numbers (the six
components of $\bC_n$ in the order $xx, yy, zz, xy, yz, xz$ and
$\theta_n$), written by the accept hook of the stepper from the converged
state and reset to $\bI$ and $\theta_0$ at the start.  The protocol is that
of Section~\ref{sec:viscoelasticity}: \code{SetLoadFactor} makes
$\Delta t = t - t_n$ the length of the step under way, between steps
$\Delta t = 0$ and the temperature rows vanish at the accepted state.  The
scheme is first order in the step; the tests measure the halving of the
error against the series solution of a conduction problem.

\subsection{Coupled weak form, kernel and module}

The unknown is $[\bu; p; \theta]$ on the Taylor--Hood spaces of
Section~\ref{sec:mixed} with the temperature on the pressure's space
(order $k - 1$; the reference's element).  The residuals of the step are
\begin{align}
  R_u(\bw) &= \int_{\Omega_R} \bigl[s(\theta)\bP_{\mathrm{iso}}(\bF) + pJ\bF^{-T}\bigr]:\Grad\bw\,\ud V + [\text{follower and contact terms}] - \sum_i s_i(t)\,[\text{dead loads}], \notag\\
  R_p(q_p) &= \int_{\Omega_R} q_p\Bigl[\frac{u'(J_m)}{J_\theta} - \frac{p}{\kappa}\Bigr]\ud V, \label{eq:thermoweak}\\
  R_\theta(q) &= \int_{\Omega_R} q\Bigl[c_v(\theta - \theta_n) - \tfrac12\theta\,\bM:(\bC - \bC_n)\Bigr]\ud V + \Delta t\int_{\Omega_R} k\,J\,\bC^{-1}\Grad\theta\cdot\Grad q\,\ud V \notag\\
  &\quad - \Delta t\sum_i s_i(t)\oint_{\Gamma_i} h_i\,|\cof\bF\,\bN|\,q\,\ud A_R, \notag
\end{align}
the last sum over the heat-flux entries of Section~\ref{sec:thermalbcs};
the prescribed temperatures are essential conditions of the third block, the
prescribed displacements of the first, and the pressure carries none.  The
kernel (\file{kernels/thermo\_mixed\_total\_lagrangian.hpp}) evaluates the
four densities $\bP$, the constraint, the capacity--heating term and the
flux $-\bQ$ of a point by one templated function of $(\bF, p, \theta, \Grad\theta)$
with the point's $\bC_n$, $\theta_n$ as parameters, and forms the nine
tangent blocks by dual seeds of that function: one seed per entry of $\bF$
that the displacement dofs reach (the in-plane block, plus the hoop stretch
of an axisymmetric point), one for $p$, one for $\theta$ and one per
component of $\Grad\theta$, each a rank-one update of the element matrices
with the B-rows of the dofs it depends on.  So $\bK_{uu}$, $\bK_{pu}$ and
$\bK_{\theta u}$ come from the $\bF$ seeds (the last through
$\bM(\bF)$, $\bC(\bF)$ and $J\bC^{-1}$), $\bK_{up}$ and $\bK_{pp}$ from the
$p$ seed, $\bK_{u\theta}$ (through $s(\theta)$), $\bK_{p\theta}$ (through
$J_\theta$) and the capacity and heating parts of $\bK_{\theta\theta}$ from
the $\theta$ seed, and the conduction part of $\bK_{\theta\theta}$ from the
gradient seeds; $\bK_{\theta p} = \bm 0$.  The tangent is the exact
derivative of \eqref{eq:thermoweak}, including $\theta\,\partial\bM/\partial\theta$
(the dual seed differentiates $\bM$ too).  The face integrators of the
mixed formulation are reused through an adapter that sizes the temperature
entries of the element vectors and matrices to zero, since the block form
keeps one set across integrators; the heat-flux face integrator fills the
temperature block and its tangent in the displacement.  The module
\code{ThermoSolidMechanicsTL} (\file{physics/thermo\_solid\_mechanics\_tl.*})
is the mixed module of Section~\ref{sec:mixed} with the third space, the
history, the thermal entries, the initial state $\theta = \theta_0$, the
nodal field \code{temperature} and the quadrature quantities with the
field pressure and the entropic modulus.  It is selected by
\code{formulation: mixed} together with the thermal block and needs the
\code{time} block; the displacement formulation with a temperature (a
penalty on the thermal expansion) and inertia with a temperature are input
errors, as is plane stress.  The linear solver is the sparse direct solver
of Section~\ref{sec:direct} on the merged $3\times3$ block Jacobian: the
saddle-point preconditioner has no temperature block.  The tangent
predictor of Section~\ref{sec:loadstepping} spreads the increments of the
prescribed displacements and temperatures through the body before Newton
starts, and the stretched cylinder of the examples needs it, since a large
prescribed increment applied to one layer of elements inverts them at the
plain start.  The energy diagnostic is the mechanical free energy
$s(\theta)\Psi_{\mathrm{iso}} + \kappa u(J_m)$ of the displacement form.

%============================================================================
\section{Scalar transport}
\label{sec:scalar}
%============================================================================

The second physics of the library is the transport of one scalar unknown
$u$ (a concentration, a temperature, a potential) by diffusion, convection
and reaction on a fixed domain: the first-order-in-time physics of the
flux/source seam, next to the solid.  It is driven by its own thin
executable, \file{apps/scalar\_transport.cpp}, whose inputs reproduce the
convection--diffusion verification drivers of \file{myapps/convection\_diffusion}
(verification manual, Section~``Scalar transport cases'').  This section
gives the strong and weak forms, the laws, the time integration, the kernel
and its tangent, the boundary conditions and the flows through them, the
error measure, and what makes a problem linear.

\subsection{Strong and weak forms}

Per unit volume of the fixed domain $\Omega$,
\begin{equation}
  c(u)\,\frac{\partial u}{\partial t} - \dive\bm{\mathcal F}(u, \grad u) + \mathcal S(u, \grad u) = 0,
  \label{eq:scalarstrong}
\end{equation}
with the flux and the source in one of two forms of the convection,
\begin{align}
  \bm{\mathcal F} &= \kappa(u)\,\grad u, & \mathcal S &= \bm\beta\cdot\grad u + s\,u - f
  && \text{(non-conservative, default)}, \notag\\
  \bm{\mathcal F} &= \kappa(u)\,\grad u - \bm\beta\,u, & \mathcal S &= s\,u - f
  && \text{(conservative)},
  \label{eq:scalarforms}
\end{align}
where $c$ is the capacity, $\kappa$ the isotropic conductivity, $\bm\beta(\bx, t)$
the velocity, $s$ the reaction coefficient and $f(\bx, t)$ the source.  The
vector $\bm{\mathcal F}$ is the negative of the physical flux (the
convention of the thermo kernel of Section~\ref{sec:thermoelasticity},
whose density is $-\bQ$), so the natural boundary condition is
$\bm{\mathcal F}\cdot\bn = 0$; the two forms coincide for a
divergence-free velocity, the non-conservative one being that of MFEM's
\code{ConvectionIntegrator} and of the myapps drivers.  The boundary carries
either the unknown itself or the inward flux $g = \bm{\mathcal F}\cdot\bn$
(for pure diffusion $\kappa\,\partial u/\partial n$, positive into the
domain, the convention of the heat-flux entries of Section~\ref{sec:thermalbcs}).

The weak form over a step $t_n \to t_{n+1} = t_n + \Delta t$ with the
implicit Euler method reads
\begin{equation}
  R(u; q) = \int_\Omega q\,c(u)\,\frac{u - u_n}{\Delta t}\ud V
          + \int_\Omega \grad q\cdot\bm{\mathcal F}(u, \grad u)\ud V
          + \int_\Omega q\,\mathcal S(u, \grad u)\ud V
          - \sum_i s_i(t)\oint_{\Gamma_i} g_i\,q\ud A = 0,
  \label{eq:scalarweak}
\end{equation}
the sum over the flux entries with their schedules $s_i$; the prescribed
values are essential conditions.  Without a \code{time} block the problem
is steady: the rate term is absent ($\Delta t = 0$) and the data follow
the schedules of the pseudo-time of Section~\ref{sec:loadstepping}, one
load step by default.

\subsection{Laws}

The capacity, the conductivity and the reaction are laws of the unknown,
each affine,
\begin{equation}
  a(u) = a_0 + a_1\,(u - u_{\mathrm{ref}}),
\end{equation}
given in the input as a number ($a_1 = 0$) or as a map
\code{\{value, slope, reference\}} (\file{kernels/materials/scalar\_transport\_model.hpp}).
This covers constant coefficients and the property laws of the nonlinear
diffusion case of the drivers; the velocity and the source are expressions
of $(\bx, t)$ evaluated at the quadrature points and never enter the
tangent.  The residual is affine in the unknown when the slopes of the
capacity and the conductivity vanish, whatever the velocity, the source and
the data do with $\bx$ and $t$: the module reports such a problem as linear
(Section~\ref{sec:scalarlinear}).

\subsection{Time integration and the accepted state}

The rate term is integrated by the implicit Euler method inside the kernel,
not by a decorator as the inertia of Section~\ref{sec:dynamics}: the
capacity may depend on the state, which puts the term in the nonlinear
form (Section~\ref{sec:seam}).  The accepted state $u_n$ is a grid function
on the space of the unknown, read at the element's dofs by the kernel and
interpolated at the quadrature points (the same values a quadrature-point
history would hold); \code{SetLoadFactor}$(t)$ makes $\Delta t = t - t_n$
the step under way and \code{AcceptStep} copies the converged state into
$u_n$, the protocol of Section~\ref{sec:viscoelasticity}.  The initial
state is the projection of the \code{initial} expression (zero without
one), evaluated at $t = 0$ whatever time the coefficient carries, and it
becomes the accepted state of $t = 0$.  Equal steps $t_{\text{final}}\,k/n$
differ in their last bits; a step within $10^{-12}$ of the one the
Jacobian was assembled for is taken as that one, so that a linear problem
keeps its operator (Section~\ref{sec:scalarlinear}).  The scheme is first
order in the step; the tests measure the rate against manufactured
solutions and against the series solution of the nonlinear diffusion
case.

\subsection{Kernel and tangent}

\code{ScalarFluxIntegrator} (\file{kernels/scalar\_flux.hpp}) is a
\code{NonlinearFormIntegrator} on the scalar $H^1$ space of order $k$.  At a
quadrature point with weight $w_q$, shape functions $N_a$ and physical
gradients $\mathsf D_{aj}$, one templated function returns the two
densities of \eqref{eq:scalarweak},
\begin{equation}
  r = c(u)\,\frac{u - u_n}{\Delta t} + \mathcal S(u, \grad u),
  \qquad
  \bm{\mathcal F}(u, \grad u),
\end{equation}
and the element residual is $R_a = \sum_q w_q\,[N_a\,r + \mathsf D_{aj}\,\mathcal F_j]$.
The tangent comes from $1 + d$ dual seeds of the same function, one on $u$
and one per component of $\grad u$; the seed $\sigma$ gives the derivative
$(r', \bm{\mathcal F}')$ of both densities, the row
$\rho_a = w_q\,[N_a\,r' + \mathsf D_{aj}\,\mathcal F'_j]$ and the
dependence $\chi_b$ of the seed on the dofs ($N_b$ for $u$, $\mathsf D_{bj}$
for $\partial u/\partial x_j$), and the element matrix accumulates the
rank-one update $K_{ab} \mathrel{+}= \rho_a\,\chi_b$: the loop of the thermo
kernel of Section~\ref{sec:thermoelasticity} for one field.  The rule is
the Gauss rule of order $2k + 3$ of Section~\ref{sec:fem}, or the order
\code{transport.quadrature\_order} when given; the source is integrated
with the same rule, which is what the cross-checks against the drivers use
to reproduce their rules.  The kernel is written for two and three space
dimensions and does not touch the solid kernels; with two instances of the
contract in hand, factoring their element loops into one CG kernel is a
later refactor.

\subsection{Boundary conditions and flows}

\code{ScalarConditions} (\file{physics/scalar\_conditions.*}) holds the
Dirichlet entries (faces by attribute, or the node nearest to a point, as
the pins of Section~\ref{sec:loads}; one expression $\bar u(\bx, t)$, a
schedule, a name) and the flux entries (the inward flux $g(\bx, t)$ per unit
area of the faces, assembled by MFEM's \code{BoundaryLFIntegrator} into one
true vector per entry, reassembled when the expression mentions $t$, and
summed with the schedules into the external vector the residual subtracts).
The default schedule is that of the solid: the ramp of the pseudo-time,
constant when the data mention $t$ or under a \code{time} block.  The flow
through a Dirichlet entry is the sum of the full residual (essential rows
kept) over its dofs,
\begin{equation}
  \Phi_i = \sum_{j \in \text{entry } i} R_j,
  \qquad
  R_j = \oint_{\partial\Omega} N_j\,\bm{\mathcal F}\cdot\bn\ud A
  \quad\text{for the exact solution},
\end{equation}
the flux entering the domain through the entry's faces: the scalar
analogue of the reactions of Section~\ref{sec:reactions}.  A corner dof
shared by two entries' faces contributes the integral over both to the entry
that owns it, as a corner node does for the reactions.

\subsection{Errors against an exact solution}

With \code{output.exact}, an expression $u_{\mathrm{ex}}(\bx, t)$, the module
computes after every accepted step the $L^2$ error with the rule of order
$2k + 3$, its ratio to $\|u_{\mathrm{ex}}\|_{L^2}$ (zero when that norm is
below $10^{-14}$) and the largest nodal error against the interpolant
$\max_i |u_i - u_{\mathrm{ex}}(\bX_i)|$, the measures of the drivers; the
executable prints them on the step line, writes them to
\file{error\_history.csv} next to the ParaView collection, and registers
the fields \code{<u>\_exact} and \code{<u>\_error}.  $H^1$ errors are
computed in the tests from coded gradients.

\subsection{Linearity, solvers and the executable}
\label{sec:scalarlinear}

When both slopes vanish the module's \code{IsLinear} is true: Newton takes
its one exact step and accepts the state at the round-off floor of the
residual (Section~\ref{sec:linearsolve}), and the Jacobian is assembled once
and reused with its stamp until the step length changes, the boundary
conditions change or the velocity expression mentions $t$ (the mechanism
of Section~\ref{sec:linearsolve}, extended by the two conditions), so that a
transient linear run with equal steps assembles one matrix and sets up one
solver for the whole history.  Nonlinear laws use the damped Newton of
Section~\ref{sec:solvers}, the tangent predictor allowed.  The linear
solver is GMRES with BoomerAMG's scalar defaults (\code{amg: scalar}, the
only mode of a space with one component, chosen automatically), CG when
there is no velocity (\code{cg\_amg} with one is an input error: the
operator is not symmetric), or the sparse direct solver of
Section~\ref{sec:direct}.

The executable \file{apps/scalar\_transport.cpp} parses the schema of
\file{base/scalar\_config.\{hpp,cpp\}} (the key \code{physics:
scalar\_transport} marks an input as its own; the solid executable rejects
that key and this one rejects the solid keys), builds the mesh and the
problem, and runs the stepper of Section~\ref{sec:loadstepping} in the
pseudo-time or in physical time with a callback that writes the ParaView
cycles, the Newton count, the errors, the flows (\code{output.flows},
\file{flows.csv}) and the probes after every accepted step; the nodal
unknown is registered under the name \code{transport.unknown} (\code{u} by
default) and the quadrature quantity \code{flux}, the vector
$\bm{\mathcal F}$, with the presentations of Section~\ref{sec:post}.

%============================================================================
\section{Boundary conditions and loading}
\label{sec:loads}
%============================================================================

All boundary conditions and the body force are entries of a \emph{load
set} (\file{physics/loads.hpp}) attached to the displacement space and
shared by both formulations.  Each entry $i$ has data (an MFEM coefficient
built from the input expressions) and a scalar schedule $s_i(t)$ of the
pseudo-time.

\subsection{Pseudo-time and schedules}

The load path is parametrised by $t \in [0, 1]$.  The data of entry $i$
enters as $s_i(t)\,f_i(\bX, t)$, where $f_i$ may itself depend on $t$
through its expression.  Schedules are piecewise linear:
\begin{equation}
  s(t) = \begin{cases}
  \text{ramp}(t_0, t_1): & 0 \ (t \le t_0),\ \ \dfrac{t - t_0}{t_1 - t_0}\ (t_0 < t < t_1),\ \ 1 \ (t \ge t_1), \\[6pt]
  \text{constant}: & 1 \ (t > 0), \\[2pt]
  \text{table}(t_k, s_k): & \text{linear interpolation of } (t_k, s_k), \text{ clamped outside}.
  \end{cases}
\end{equation}
The default is the ramp over $[0, 1]$, unless the expression mentions $t$,
in which case it is the constant schedule and $t$ enters through the
function only; a schedule given together with such an expression
multiplies it.  With the ramp on every entry, $t < 1$ solves a
proportionally scaled problem and $t = 1$ the problem of the weak form.
Since hyperelasticity without follower loads is conservative, the converged
state at $t = 1$ does not depend on the path; staged loading is a test of
the implementation, not a different problem.  Follower pressures are not
conservative and can produce path dependence in the presence of
instabilities.

\paragraph{Physical time.}
In a dynamic analysis (Section~\ref{sec:dynamics}) $t$ is the physical time
in $[0, t_{\text{final}}]$.  Ramps and tables are given on that interval (a
ramp without bounds spans it), the value at $t = 0$ is the right limit, so
the constant schedule is $1$ for $t \ge 0$, and an entry \emph{without} a
schedule is constant: its data is the expression as it stands, a step load
when the expression does not mention $t$.  The quasi-static default ramp
over the pseudo-time has no meaning there.  Because that default depends on
the analysis, the executable prints the resolved time dependence of every
entry in the header of a dynamic run.

\subsection{Prescribed displacements}

A Dirichlet entry prescribes, on a set of boundary attributes, either all
components or the listed subset (\code{components}), which is how rollers
and symmetry planes are expressed.  The essential true degrees of freedom
of each entry are collected per component; their union is the essential
set of the problem.  Before each Newton solve the stepper overwrites the
essential entries of the solution vector with $s_i(t)$ times the boundary
projection of $\bar\bu_i$ (\code{ApplyDirichlet}); the residual is zeroed
and the Jacobian rows and columns eliminated there.  Two entries may
overlap on shared edges or corners; the later entry wins and a warning
names entries that prescribe the same component on the same attribute.

\paragraph{Point constraints.}
An entry with \code{point} in place of the attributes prescribes its
components at the displacement node nearest to the given point: the owning
rank is found by a global minimum of the distance over the true dofs, the
entry's essential dofs are that node's components, and the data is the
coefficient's value at the node (the projection of $\bar\bu_i$ read there).
A point farther than $10^{-8}$ of the mesh diameter from every node is an
input error.  Nothing else changes: the schedule, the elimination and the
reaction (the nodal force at the node) are those of a face entry.  This is
how single nodes are pinned to remove rigid modes, as the reference codes
of Section~\ref{sec:thermoelasticity} do; an edge of nodes takes one entry
per node.

\subsection{Dead loads}

Nominal tractions $\bar\bT$ (per unit reference area), dead normal
pressures $\bar\bT = -p\bN$ on the reference outward normal, and the body
force $\rho_R\bb$ are assembled once each into true-dof vectors $L_i$ with
stock MFEM linear-form integrators; entries whose expression depends on
$t$ are reassembled whenever $t$ changes.  The external load at time $t$
is $\sum_i s_i(t)\,L_i$ and is subtracted from the internal force in the
residual.

\subsection{Follower pressure}
\label{sec:follower}

A pressure $p$ per unit \emph{current} area on $\Gamma_R^p$ exerts the
spatial traction $-p\bn$; by Nanson's formula \eqref{eq:nanson} its
nominal counterpart is $\bar\bT(\bu) = -p\,J\bF^{-T}\bN = -p\,\cof\bF\,\bN$,
which depends on the displacement.  It is therefore a boundary term of the
nonlinear form,
\begin{equation}
  R_{\text{fol}}(\bu; \bw) = \sum_i s_i(t)\int_{\Gamma_R^{p_i}} p_i\,(\cof\bF\,\bN)\cdot\bw\ud A ,
\end{equation}
with the load stiffness
\begin{equation}
  D_{\Delta\bu}\int_{\Gamma_R^p} p\,(\cof\bF\,\bN)\cdot\bw\ud A
  = \int_{\Gamma_R^p} p\,\bigl[(\bF^{-T}:\Grad\Delta\bu)\,J\bF^{-T}\bN
    - J\bF^{-T}(\Grad\Delta\bu)^T\bF^{-T}\bN\bigr]\cdot\bw\ud A ,
  \label{eq:foltangent}
\end{equation}
using $DJ = J\bF^{-T}:\Grad\Delta\bu$ and
$D\bF^{-T} = -\bF^{-T}(\Grad\Delta\bu)^T\bF^{-T}$.  This tangent is not
symmetric, so the conjugate gradient option of the linear solver is refused
for inputs with a follower pressure.  The dead pressure and the follower
pressure coincide at $\bF = \bI$ and differ at $O(p^2)$.  The
implementation (\file{kernels/follower\_pressure.hpp}) is described in
Section~\ref{sec:folkernel}.

\subsection{Rigid-sphere penalty contact}
\label{sec:contact}

A \code{contact} entry of the boundary conditions (\code{type:
rigid\_sphere}) presses a rigid sphere of radius $r$ and centre $\bm c(t)$,
an expression of the time, against a boundary $\Gamma_R^c$ through the
penalty of the reference's codes: per unit reference area the energy
\begin{equation}
  E = \frac{k}{2}\,\bigl\langle r^2 - |\bx - \bm c|^2\bigr\rangle_+^2, \qquad \bx = \bX + \bu,
\end{equation}
whose derivative is the traction on the body
\begin{equation}
  \bt = -\frac{\partial E}{\partial\bu} = 2k\,g\,(\bx - \bm c), \qquad
  g = \bigl\langle r^2 - |\bx - \bm c|^2\bigr\rangle_+ ,
\end{equation}
which pushes whatever lies inside the sphere out of it along the radius
and vanishes outside.  The boundary term of the nonlinear form and its
tangent are
\begin{equation}
  R_{\text{con}}(\bu; \bw) = -s(t)\int_{\Gamma_R^c}\bt\cdot\bw\ud A, \qquad
  \frac{\partial\bt}{\partial\bu} = 2k\bigl[g\,\bI - 2(\bx - \bm c)\otimes(\bx - \bm c)\bigr] \quad (g > 0),
\end{equation}
symmetric, with the kink of the penalty at $g = 0$; $s(t)$ is the entry's
schedule on the penalty (constant by default; the indenter moves through
$\bm c(t)$).  The resultant force and moment of $\bt$ on the body are
reported with the reactions under the entry's name: the physics keeps a
second form holding the contact integrator alone, whose residual at the
accepted state gives the nodal contact forces, summed like a reaction
(Section~\ref{sec:reactions}).  The face kernel
(\file{kernels/rigid\_sphere\_contact.hpp}, both formulations) integrates
with the rule of the follower pressure.

\subsection{Thermal boundary conditions}
\label{sec:thermalbcs}

The coupled problem of Section~\ref{sec:thermoelasticity} adds two kinds of
entry on the temperature space, kept by that module rather than by the load
set.  A \code{temperature} entry prescribes $\theta = s_i(t)\bar\theta_i$ on
the faces of its attributes (one expression; under the \code{time} block an
entry without a schedule is constant, so the expression itself is the
history): its essential dofs of the temperature block are treated as those
of the displacement.  A \code{heat\_flux} entry applies an inward heat flux
$h$, by default per unit current area, which in the reference configuration
is $h\,|\cof\bF\,\bN|$ per unit reference area (Nanson), so that the term
depends on the displacement and carries a tangent in it by dual numbers over
the face dofs, as the follower pressure; with \code{per\_unit:
reference\_area} it is $h$ itself.  The entry contributes
$-\Delta t\,s_i(t)\oint h\,|\cof\bF\,\bN|\,q\,\ud A_R$ to the temperature
residual of the step.  Insulated is the natural condition; convection,
radiation, volumetric sources and thermal contact are not provided.

\subsection{Reactions}
\label{sec:reactions}

The reaction of a Dirichlet entry is the resultant of the nodal forces the
constraint exerts on the body: the residual \eqref{eq:weak} evaluated with
the essential rows \emph{not} zeroed, restricted to the entry's essential
degrees of freedom (internal minus external nodal forces).  This is the
exact discrete counterpart of the traction integral over the constrained
face.  The moment is taken about the origin at the current nodal positions
$\bX + \bu$.  Two entries that share nodes both count the shared forces.
The reactions are printed after every step and, with ParaView output,
written to \file{<paraview>/reactions.csv}, one row per accepted step from
the initial state on (\code{step, t}, then \code{fx, fy, fz, mx, my, mz}
per entry), which lines up with the cycles of the \file{.pvd} when
\code{output.every} is 1.  Post-processing should use these rather than a
stress field integrated over the face: the nodal stress of the output or
the finite element stress at quadrature points is off by up to $10\%$
where the face meets free faces or clamped corners (verification manual,
the finite elasticity examples).

\subsection{Expressions}

Boundary data and the body force are strings $f(x, y, z, t)$ in the
reference coordinates and the pseudo-time ($z = 0$ in two dimensions),
parsed once into a postfix program and evaluated per point without
allocation.  The grammar has numbers, \code{x y z t}, \code{pi}, the
operators \code{+ - * / \^{}} (\code{\^{}} binds tightest, associates to the
right, and unary minus binds looser, so \code{-x\^{}2} $= -(x^2)$),
parentheses, the functions \code{sin cos tan exp log sqrt abs pow min max},
and \code{if(cond, a, b)} with the comparisons \code{< <= > >= == !=}
evaluating to 1 or 0.  Expressions cannot refer to the solution; a load
that depends on the deformation is a follower load, and only the follower
pressure is provided.

%============================================================================
\section{Finite element discretisation}
\label{sec:fem}
%============================================================================

\subsection{Spaces and quadrature}

The displacement lives in the vector $H^1$ space of polynomial order $k$
(\code{mesh.order}) on the reference mesh, ordered by vector dimension
(the components of a node are consecutive), so that true-dof vectors can be
read node by node; the pressure of the mixed formulation in the scalar
$H^1$ space of order $k-1$.  Quadrilaterals, triangles, hexahedra and
tetrahedra, including mixed and curved (second-order geometry) meshes, are
accepted.  Volume and face integrals use the Gauss rule of order $2k + 3$
on the element geometry; the same rule defines the quadrature space of the
post-processed quantities (Section~\ref{sec:post}).

\subsection{Element residual and tangent of the displacement formulation}

At a quadrature point with weight $w_q$ (the rule weight times the
reference map's Jacobian determinant) the physical shape-function gradients
$\mathsf{D}_{aj} = \partial N_a/\partial X_j$ give
$H_{ij} = \sum_a u_{ai}\mathsf{D}_{aj}$ and the material returns $\bP$ and
$\bA$ at $\bF = \bI + \bH$.  The element residual and tangent are
\begin{equation}
  r_{ai} = \sum_q w_q\,P_{ij}\,\mathsf{D}_{aj},
  \qquad
  K_{(ai),(bk)} = \sum_q w_q\,\mathsf{D}_{aj}\,A_{ijkl}\,\mathsf{D}_{bl},
  \label{eq:element}
\end{equation}
summed over the in-plane indices $j, l < d$ (in two dimensions the
out-of-plane components of $\bA$ are never used).  The loops are written in
\file{kernels/total\_lagrangian.hpp} as plain host code without allocation
or virtual calls; the quadrature-point bodies (deformation gradient,
stress, tangent, Cauchy stress) are free functions on plain tensors.  The
integrator is a template over the material type, instantiated once per
model through \code{std::visit} at setup, so no dispatch happens per
quadrature point.  The energy $\int\Psi\ud V$ is provided by the same
integrator when the material has an energy, and is used only as a
diagnostic.

\subsection{Element blocks of the mixed formulation}

The block integrator (\file{kernels/mixed\_total\_lagrangian.hpp}) adds the
pressure shape functions $\mathsf{N}_b$ and forms, per quadrature point,
\begin{equation}
  \begin{aligned}
  r^u_{ai} &= w_q\,\bigl[P_{\mathrm{iso},ij} + p\,(J F^{-T})_{ij}\bigr]\mathsf{D}_{aj}, &
  r^p_b &= w_q\,\Bigl(u'(J) - \frac{p}{\kappa}\Bigr)\mathsf{N}_b, \\
  K^{uu}_{(ai),(bk)} &= w_q\,\mathsf{D}_{aj}A_{ijkl}\mathsf{D}_{bl}, &
  K^{up}_{(ai),b} &= w_q\,(JF^{-T})_{ij}\mathsf{D}_{aj}\mathsf{N}_b, \\
  K^{pu}_{b,(ai)} &= u''(J)\,K^{up}_{(ai),b}, &
  K^{pp}_{ab} &= -\frac{w_q}{\kappa}\mathsf{N}_a\mathsf{N}_b,
  \end{aligned}
\end{equation}
with $p = \sum_b p_b\mathsf{N}_b$, $\bA$ the dual tangent of
$\bP_{\mathrm{iso}} + pJ\bF^{-T}$ at fixed $p$, and $u', u''$ the
normalised derivatives of the volumetric law.  In two dimensions $J$ and
$J\bF^{-T}$ are those of the $2\times2$ block (plane strain).  The energy
of this integrator is the mixed functional
$\Psi_{\mathrm{iso}} + p(J-1) - \kappa\,u^*(p/\kappa)$, with $u^*$ the
Legendre transform of the volumetric law
(Section~\ref{sec:post}); the last term is absent when $\kappa = \infty$.

\subsection{Follower-pressure face kernel}
\label{sec:folkernel}

On each boundary face of $\Gamma_R^p$ the integrator takes the adjacent
element's displacement degrees of freedom, evaluates $\bH$ and
$\bF = \bI + \bH$ at the face quadrature points, and obtains $\bN\ud A$ as
the area-weighted reference normal of the face map (\code{CalcOrtho}),
oriented outward by a one-time check against the vector from the element
centre to the face centre.  The face residual is
$r_{ai} = \sum_q s\,\hat w_q\,p\,N_a\,(\cof\bF\,\bN)_i$ with $\hat w_q$ the
face rule weight; the tangent is obtained by seeding $\bF$ with dual numbers
in the direction of each element displacement degree of freedom
($\Delta u_k = N_b \Rightarrow \ud F_{kj} = \mathsf{D}_{bj}$) and reading
the derivative part of $\cof\bF\,\bN$, which reproduces
\eqref{eq:foltangent} without writing it out.  The scale $s$ is the
schedule value $s_i(t)$, owned by the load set and read through a pointer
so that the integrator needs no rebuild when $t$ changes.  In the mixed
formulation the same kernel acts on the displacement block only.

\subsection{The assembled operator}

Each physics module is an \code{mfem::Operator} on true degrees of
freedom: \code{Mult} returns the residual \eqref{eq:weak} or
\eqref{eq:mixedweak} with the essential rows zeroed;
\code{GetGradient} returns the assembled \code{HypreParMatrix} (or the
$2\times2$ \code{BlockOperator} of such matrices) with the essential rows
and columns eliminated.  Full assembly through
\code{ParNonlinearForm} and \code{ParBlockNonlinearForm} is used; partial
assembly is not implemented.

%============================================================================
\section{Automatic differentiation of the tangents}
\label{sec:ad}
%============================================================================

A dual number $a = (a_v, a_d)$ carries a value and one directional
derivative with the arithmetic $(a + b) = (a_v + b_v, a_d + b_d)$,
$(ab) = (a_vb_v, a_db_v + a_vb_d)$, and the chain rule for $/$, $\log$,
$\exp$, $\sqrt{\ }$, $\mathrm{pow}$, $|\cdot|$, $\sin$, $\cos$
(\file{base/dual.hpp}).  The fixed-size tensors of \file{base/tensor.hpp}
are templated on the scalar type, so any material written with their
operations compiles for $T = $ \code{dual} unchanged.

The material tangent \eqref{eq:tangent} is obtained by seeding: for each
in-plane index pair $(k, l)$ the entry $F_{kl}$ is given derivative part
one, $\bP$ is evaluated once with duals, and its derivative parts fill
$A_{ijkl}$.  This costs $d^2$ stress evaluations per quadrature point (nine
in three dimensions, four in plane strain) and produces the exact tangent
of the implemented stress to round-off, with no hand-coded derivatives.
The same seeding gives the tangent of $\bP_{\mathrm{iso}} + pJ\bF^{-T}$ at
fixed $p$ in the mixed kernel and the follower load stiffness.  Two places
need more than plain forward mode: the Ogden model, whose eigenvectors are
not differentiable at coincident stretches (Section~\ref{sec:ogden}), and
the compressible plane-stress adapter, whose thickness stretch is the root
of a scalar equation (Section~\ref{sec:planestress}); both supply the
derivative part analytically or implicitly inside their dual overloads,
so the seeding loop is unchanged.  All tangents are checked against central
finite differences in the tests.

%============================================================================
\section{Solution algorithms}
\label{sec:solvers}
%============================================================================

\subsection{Quasi-static load stepping}
\label{sec:loadstepping}

The stepper (\file{solvers/quasi\_static.hpp}) advances $t$ through a
sequence of breakpoints $0 < t_1 < \dots < t_n = 1$ (equal increments from
\code{load\_steps}, or the segments of \code{steps}).  For each increment
it sets the pseudo-time of the load set (schedules evaluated, coefficients'
time set, time-dependent dead loads reassembled, follower scales updated),
applies the Dirichlet data to the current solution vector, and runs Newton
warm-started from the last converged state.  If Newton fails and
\code{substep.on\_failure} is set, the last converged state is restored,
the increment is halved and retried, at most \code{max\_bisections} times
and never below \code{min\_dt}; after a converged reduced increment the
stepper aims at the planned breakpoint again and never grows beyond the
planned grid.  Each accepted increment triggers the output callback
(fields, ParaView cycle, probes, reactions and their CSV).

\paragraph{Time stepping.}
A dynamic analysis is advanced by the same loop (\code{SolveDynamic}) over
targets $0 < t_1 < \dots < t_n = t_{\text{final}}$ from \code{dynamics.dt}
or the segments of \code{dynamics.steps}.  The problem is then the step
equation \eqref{eq:step}: setting the time begins a step (increment,
$c_M$, $\bm u^*$, $\bm h_n$, the time of the loads), and an accept hook,
called once per accepted increment, advances the history
$(\bm u_n, \bm v_n, \bm a_n, \bm S_n)$.  A rejected increment leaves the
history untouched, so bisection needs nothing else; a run with a halved
step equals the run with the two half steps planned.  Two details differ
from the pseudo-time.  A target counts as reached within $10^{-9}$ of the
planned increment, not within an absolute $10^{-14}$, which is below the
spacing of floating-point numbers once $t > 100$ and would leave a
degenerate last step with $c_M \sim \Delta t^{-2}$.  And Newton starts from
the last converged state with the new Dirichlet data, never from an
extrapolation in time: a better start makes $\|R_0\|$ small, and
$\text{rtol}\,\|R_0\|$ then lies below the round-off floor of the residual,
the trap that excludes the tangent predictor for linear problems.
\code{substep.min\_dt} is a time in a dynamic analysis ($10^{-3}$ of the
smallest planned step by default); \code{load\_steps}, \code{steps} and
\code{predictor} of the solver section are input errors.

\paragraph{Quasi-static analysis in physical time.}
A \code{time} block (\code{t\_final} and \code{dt} or \code{steps}, the
keys of the dynamics block without a scheme) runs the quasi-static problem
through the same loop in physical time (\code{SolveInTime}, of which the
dynamic analysis is the same function under its older name): the loads
follow the physical time exactly as under \code{dynamics} (schedules on
$[0, t_{\text{final}}]$, an entry without one is constant), the increments
are the time steps, and the accept hook advances a rate-dependent
material's history by the step length (Section~\ref{sec:viscoelasticity}).
Nothing else changes: there is no inertia, Newton and bisection are those
of the pseudo-time, the tangent predictor stays available (the start of a
step is not an extrapolation in time), and \code{substep.min\_dt} is a
time.  \code{load\_steps} and \code{steps} of the solver section are input
errors with the block, and the two blocks exclude each other.

\paragraph{Initial guess of an increment.}
By default the increment starts from the last converged state $x_n$ with
the new Dirichlet values written on the boundary.  The interior lags, so
the first Newton linearisation is about a state whose boundary layer of
elements carries the whole increment; in a sheared block the iterate that
follows can have $J$ between $0.4$ and $3$ at quadrature points, and the
next step can invert elements.  With \code{solver.predictor: tangent} the
stepper instead imposes the increment on a linear solve about the converged
state.  With $g$ the state $x_n$ carrying the new boundary values and
$d = g - x_n$ (nonzero on the essential degrees of freedom only),
\begin{equation}
  J(x_n)\,s = R(x_n) + J(x_n)\,d, \qquad x^{(0)} = x_n + d - s ,
\end{equation}
where $R$ is evaluated with the loads at the new pseudo-time and $s$
vanishes on the essential rows.  Because the assembled Jacobian has its
essential columns eliminated, $J(x_n)d$ is formed by a forward difference of
the residual, $[R(x_n + \epsilon d) - R(x_n)]/\epsilon$.  This is the start
that codes imposing the boundary condition on the Newton update obtain
implicitly.  The converged states are unchanged; Newton typically needs
half the iterations and no damped steps.  The predictor matters most for
nearly incompressible materials and is needed in practice for the
logarithmic volumetric law, whose pressure $\ln J/J$ is not monotone beyond
$J = e$ and whose stiffness $u''$ varies by two orders of magnitude over
the range of $J$ of a lagging-interior iterate.  It is a no-op for
traction-driven increments, and it falls back to the default start if the
predicted state has a non-finite residual.

\subsection{Damped Newton}

For $R(x) = 0$ with Jacobian $J(x) = \partial R/\partial x$
(\file{solvers/newton.hpp}):
\begin{enumerate}[nosep]
  \item evaluate $R_0 = R(x_0)$; converged if $\|R\| \le \text{atol}$ or
    $\|R\| \le \text{rtol}\,\|R_0\|$, where $\|\cdot\|$ is the global
    $\ell_2$ norm over the communicator and $R_0$ is the residual at the
    start of the increment (after the Dirichlet data are applied);
  \item solve $J\,\delta = R$ with the linear solver; a Krylov solve that
    hits its iteration cap is reported and the inexact step is used;
  \item Armijo backtracking: with $\alpha = 1$, accept $x - \alpha\delta$
    if $\|R(x - \alpha\delta)\| \le (1 - c\,\alpha)\|R(x)\|$ and the norm is
    finite, otherwise halve $\alpha$, at most \code{max\_halvings} times
    ($c = 10^{-4}$ by default).  A non-finite trial residual (an inverted
    element, a Gent state beyond locking) is simply rejected;
  \item stop with failure when the line search cannot reduce the residual
    or after \code{max\_it} iterations; a problem that declares itself
    linear is instead accepted at the round-off floor of its residual
    (Section~\ref{sec:linearsolve}).
\end{enumerate}
The report carries the iteration history and the failure reason, which the
stepper uses to decide on bisection.

\subsection{Linear solver of the displacement formulation}

The Jacobian is a \code{HypreParMatrix}; the solver
(\file{solvers/linear\_solver.hpp}) is GMRES (default, restart
\code{krylov\_dim}) or CG, preconditioned by BoomerAMG with either the
\emph{elasticity} options (rigid-body-mode interpolation, needs the finite
element space) or the plain \emph{systems} options (unknown-based
coarsening for $d$ components).  MFEM's Krylov tolerances are relative to
the left-preconditioned residual $\|M^{-1}r\|$, which with AMG behaves like
an error norm, so the true residual reduction is typically 10 to 100 times
weaker than \code{rtol}.  The elasticity options converge quickly in two
dimensions but stall on slender three-dimensional $k = 2$ meshes; they are
given at most 100 iterations, after which the solver rebuilds the
preconditioner with the systems options, repeats the solve, and keeps the
systems options for the rest of the run.  The AMG hierarchy is rebuilt for
every Newton Jacobian, except for a linear problem, whose one Jacobian and
hierarchy serve the whole load path (Section~\ref{sec:linearsolve}).

\subsection{Saddle-point solver of the mixed formulation}
\label{sec:saddlesolver}

Each Newton system has the structure \eqref{eq:saddle},
$\mathcal J = \bigl[\begin{smallmatrix}\bK & \bBm\\ \bBm_{pu} & -\bC\end{smallmatrix}\bigr]$
with $\bC = \bM_p/\kappa$ (zero when incompressible) and
$\bBm_{pu} = \bBm^T$ for the quadratic volumetric law only.  When the pressure is
comparable to the shear modulus the raw $\bK$ carries an indefinite
$p$-weighted geometric stiffness, which AMG handles poorly.  The solver
(\file{solvers/saddle\_point\_solver.hpp}) therefore first rewrites the
system in the algebraically equivalent \emph{augmented Lagrangian} form:
with $\gamma = \code{augmentation}\times\mu$ and
$\bm W = \diag(\bM_p)^{-1}$, adding $\gamma\bBm\bm W$ times the second block
row to the first gives
\begin{equation}
  \begin{bmatrix} \bK + \gamma\bBm\bm W\bBm_{pu} & \bBm - \gamma\bBm\bm W\bC \\ \bBm_{pu} & -\bC \end{bmatrix}
  \begin{bmatrix} \Delta\bu \\ \Delta p \end{bmatrix}
  = \begin{bmatrix} b_u + \gamma\bBm\bm W b_p \\ b_p \end{bmatrix},
  \label{eq:augmented}
\end{equation}
whose displacement block stays well conditioned (the augmentation acts as
a penalty of strength $\gamma$ on the discrete constraint).  This is solved
by flexible GMRES with the block upper-triangular preconditioner
\begin{equation}
  \Delta p = -\tilde{\bS}^{-1} r_p, \qquad
  \Delta\bu = \tilde{\bK}^{-1}\bigl(r_u - \tilde{\bBm}\Delta p\bigr),
\end{equation}
where $\tilde{\bK}^{-1}$ is an inner GMRES + BoomerAMG (systems options)
solve of the augmented displacement block to \code{inner\_rtol}, and
$\tilde{\bS}$ approximates the Schur complement by a scaled pressure mass
matrix,
\begin{equation}
  \tilde{\bS} = \frac{1 + \mu/\kappa}{\mu + \gamma}\,\bM_p ,
\end{equation}
applied by CG with Jacobi smoothing.  The scale interpolates between the
classical estimates: without augmentation it is
$\bM_p(1/\mu + 1/\kappa)$, the sum of the incompressible Stokes-like
estimate $\bBm^T\bK^{-1}\bBm \approx \bM_p/\mu$ and the penalty term
$\bM_p/\kappa$; in the incompressible limit it is $\bM_p/(\mu + \gamma)$,
the augmented-Lagrangian estimate.  On the Cook and manufactured-solution
problems this takes about 20 outer iterations independent of the mesh.
\code{augmentation: 0} disables the transformation.  The shear modulus that
scales $\gamma$ and $\tilde{\bS}$ is the base material's small-strain value
(Section~\ref{sec:moduli}).  The row operation multiplies by $\bBm\bm W$,
the $(1,2)$ block itself, not by the transpose of the $(2,1)$ block; the
two coincide only for the quadratic law, and the distinction is what keeps
Newton quadratic with the other laws.

\paragraph{Dynamic analysis.}
The displacement block is then $\bK + c_M\bM$, and for small time steps the
mass dominates it.  The Schur complement of a mass-dominated block is a
discrete pressure Laplacian, $\bBm_{pu}(c_M\bM)^{-1}\bBm$, which a scaled
pressure mass does not represent: on the quarter annulus of the tube
problem the outer iterations rise from $18$ (quasi-static) to $43$ at
$\Delta t = 0.002$ and, one refinement on, to $81$.  Following Cahouet and
Chabard~\cite{cahouetchabard} the two limits are added,
\begin{equation}
  \tilde{\bS}^{-1} = \Bigl[\frac{1 + \mu/\kappa}{\mu + \gamma}\bM_p\Bigr]^{-1}
  + \Bigl[\frac{1}{c_M}\bBm_{pu}\bm D^{-1}\bBm + \bC\Bigr]^{-1},
\end{equation}
with $\bm D$ the lumped mass (the diagonal of $\bM$ scaled to the total
mass, positive on every element type and order, which row sums are not) and
the second inverse by CG with BoomerAMG.  The second matrix is formed from
the blocks of the Jacobian at hand, so it carries their boundary conditions
(prescribed displacements: Neumann; free surfaces: Dirichlet) without a
pressure Laplacian being assembled.  The count becomes $6$ to $16$ over
three decades of $\Delta t$, independent of the mesh.  The term exists only
when the dynamic problem hands the solver $c_M$.

\subsection{Sparse direct solver}
\label{sec:direct}
\label{sec:directsolver}

With \code{solver.linear.type: direct} every Newton system is solved by a
sparse LU factorization instead (MUMPS through PETSc,
\file{src/solvers/direct\_solver.hpp}), in either formulation and in
parallel.  The Jacobian of the mixed formulation is merged into one matrix,
and the zero $(p,p)$ block of an incompressible material is left to the
pivoting of the factorization.  The matrix is refactored at every Newton
iteration, except for a linear problem
(Section~\ref{sec:linearsolve}); the Krylov and AMG keys are unused.  The
nested Krylov solvers above are built for problems too large to factor; on
small ones they are the wrong tool.  The uniaxial cube of the verification
manual's coupled-theories set ($2312$ unknowns, $100$ load steps) spends
$84\%$ of its instructions in the saddle-point solve, about $250$ AMG cycles
per load step, and runs in $433$\,s; with the direct solver it runs in
$23$\,s with the same reaction to $13$ digits, and the twisted column
($59\,028$ unknowns) goes from about five minutes to eight seconds per load
step.  Memory and factorization time grow faster than linearly with the
size, in 3D roughly as $n^{4/3}$ and $n^{2}$, so beyond some $10^5$ to
$10^6$ unknowns the iterative solvers take over.  The factorization uses one
thread per MPI rank.

%============================================================================
\section{Post-processing}
\label{sec:post}
%============================================================================

Two kinds of output fields exist (\file{physics/quadrature\_fields.hpp}).
The nodal unknowns, \code{displacement} and (mixed) \code{pressure}, are
the $H^1$ fields the solver computes and are written as they are.  The
quadrature quantities are computed once per load step at the quadrature
points of the kernels' rule into a \code{QuadratureFunction}, the source of
truth, from the state $(\bF, \bP, \Psi)$ at each point and the measures of
the material's kinematics formed from it (finite strain below; for a
small-strain material $\bsig = \bP$, the volume ratio is $1 + \tr\beps$ and
the strain is $\beps$, Section~\ref{sec:smallstrain}):
\begin{center}
\small
\begin{tabular}{@{}l l l@{}}
\toprule
Name & Components & Definition \\
\midrule
\code{cauchy\_stress} & 6 ($xx\,yy\,zz\,xy\,yz\,xz$) & $\bsig = J^{-1}\bP\bF^T$, symmetrised \\
\code{pk1\_stress} & 9 (row-major) & $\bP$ \\
\code{deformation\_gradient} & 9 (row-major) & $\bF$ (completed $\bF_c$ under plane stress) \\
\code{strain} & 6 ($xx\,yy\,zz\,xy\,yz\,xz$) & Green--Lagrange $\bE = \tfrac12(\bF^T\bF - \bI)$ \\
\code{jacobian} & 1 & $J = \det\bF$ \\
\code{vonmises} & 1 & $\sigma_{vm} = \sqrt{\tfrac32\,\bsig':\bsig'}$, $\bsig' = \dev\sym\bsig$ \\
\code{energy\_density} & 1 & $\Psi(\bF)$; mixed: $\Psi_{\mathrm{iso}} + p(J-1) - \kappa\,u^*(p/\kappa)$ \\
\code{thickness\_stretch} & 1 & $F_{33} = \lambda_3$ (plane stress only) \\
\bottomrule
\end{tabular}
\end{center}
In two dimensions $\bF$ carries $F_{33}$, so the out-of-plane terms are
included.  Derived scalars are formed at the quadrature points and then
presented, never computed from a presented field.  Each quantity is
presented where \code{output.quadrature\_at} asks:
\begin{description}[nosep]
  \item[\code{nodes}] a continuous $H^1$ field of the mesh order.
    \code{nodal\_projection: averaged} projects the data onto each
    element's own polynomial space (element mass matrix
    $\bM_e a = \int N f$) and takes the arithmetic mean of the element
    values at shared nodes (the classical extrapolate-and-average,
    communicated across ranks); \code{projected} solves the global $L^2$
    projection with the consistent mass matrix (CG to $10^{-14}$), which is
    optimal in the least-squares sense but can overshoot near stress
    concentrations.
  \item[\code{elements}] the quadrature-weighted volume average per element,
    an order-0 field named \code{<name>\_elem}.
  \item[\code{quadrature\_points}] the raw values, written as a point cloud
    \code{<name>\_qp} (one VTU per rank, listed in the same PVD as the
    mesh).
\end{description}
For a homogeneous state all presentations agree to machine precision; on a
real problem their differences are discretisation error.  Probes evaluate
every registered grid function (nodal unknowns, \code{<name>} and
\code{<name>\_elem}) at physical points, averaging when several ranks find
the point.  The internal energy $\int\Psi\ud V$ (or the mixed functional)
is printed at the end; it is computed by a domain-only form, so follower
pressure work is not included (that load is not conservative).  For the
mixed formulation the integrand is
$\Psi_{\mathrm{iso}} + p(J-1) - \kappa\,u^*(p/\kappa)$, where
$u^*(\pi) = \sup_J[\pi(J-1) - u(J)]$ is the Legendre transform of the
normalised volumetric law ($\pi^2/2$ for the quadratic law, so that the
integrand is the perturbed Lagrangian \eqref{eq:lagrangian}; closed forms
for the Simo--Taylor and $J\ln J$ laws, a scalar Newton solve for the
logarithmic one).  The stored energy $\Psi_{\mathrm{iso}} + U(J)$ is
deliberately not used: the discrete displacement satisfies the constraint
only weakly, and $\kappa$ times the square of its pointwise violation
would dominate the value for nearly incompressible materials.  The last
term is absent when $\kappa = \infty$, where the integrand is the
Lagrangian and equals $\Psi_{\mathrm{iso}}$ on the constraint.

In a dynamic analysis the nodal fields \code{velocity} and
\code{acceleration} of the time integrator join the unknowns (written and
probed like the displacement), \code{output.energy} prints after every step
the kinetic and the internal energy, the external work and their balance
(Section~\ref{sec:dynamics}), and \code{output.every} writes every $n$-th
step to ParaView (and always the last); the probe, reaction and energy
lines stay per step, which is how time histories are read from a log.

%============================================================================
\section{Mesh input}
%============================================================================

Meshes come from files (Gmsh 2.2 or 4.x, or any format MFEM reads).  Every
physical group becomes an element or boundary attribute numbered by its tag,
and the group names are kept as attribute sets, so boundary conditions and
material regions may name groups instead of numbers; both are checked
against the mesh and the error lists what the mesh provides.  The loader
constructs the mesh in place because MFEM's move operations drop the
attribute sets.  Options: \code{serial\_refine} and \code{parallel\_refine}
uniform refinements (a second-order Gmsh mesh refines along the curved
geometry, independently of the finite element order), and \code{perturb},
a deterministic pseudo-random jitter of the interior vertices by at most
the given fraction of the smallest element size in each coordinate
(straight-sided meshes only, applied before refinement so that all levels
share one distortion pattern; inverted elements are refused).  The
Cartesian box and the bilinear corner map remain available
programmatically for the tests.

%============================================================================
\section{Framework architecture and state of the seam}
\label{sec:seam}
%============================================================================

The architecture document defines a \emph{flux/source seam}: a physics
declares a flux $\bm{\mathcal F}(u, \nabla u)$, a source $\mathcal S(u)$ and
a numerical flux $\hat{\bm{\mathcal F}}(u^-, u^+, \bn)$ per quadrature
point, and the framework turns one definition into CG volume, DG face or
HDG trace operators, with partial or full assembly and device execution.
Solid mechanics fits this shape with $\bm{\mathcal F} = \bP(\bF)$ in total
Lagrangian form.  The present code implements the seam as the minimum its two physics need:
the solid (three element integrators) and the scalar transport
(Section~\ref{sec:scalar}, the first instance of a scalar flux/source
kernel with a source term, a capacity and a flux condition).  The following
is still specific to each and would generalise into the framework's
kernels:

\begin{itemize}[nosep]
  \item The displacement integrator hard-codes the flux $\bP(\bF)$
    contracted with $\Grad\bw$, a vector $H^1$ unknown and the contraction
    $\Grad\bw:\bA:\Grad\Delta\bu$.  The quadrature-point bodies are free
    functions on plain tensors and can be lifted into a general
    $\bm{\mathcal F}(u, \nabla u)$, $\mathcal S(u)$ contract; the element
    loops would become the framework's CG volume kernel.
  \item The mixed kernel is a second block integrator with its own loops,
    and the scalar flux kernel a third, single-field one whose contract
    ($\bm{\mathcal F}(u, \grad u)$, $\mathcal S(u, \grad u)$, a capacity)
    is the seam's for one unknown; a multi-field CG kernel over a list of
    unknowns with a block flux/source contract would absorb all three.  The
    saddle-point solver is specific to the $2\times2$ $u$--$p$ structure and
    the pressure-mass Schur approximation.
  \item Sources are not part of the solid integrators: the body force is a
    dead load on the linear-form side.  The scalar kernel carries its
    source and its reaction inside the nonlinear form with their tangent,
    the shape a reaction or heat source of the solid would take.
  \item Boundary terms are stock linear-form integrators plus one
    hand-written face integrator (the follower pressure).  Robin and flux
    conditions and DG/HDG numerical fluxes have no home yet.
  \item Dirichlet conditions of the solid prescribe components of the
    vector unknown on an attribute (\code{LoadSet}); those of the scalar
    unknown live in a second class (\code{ScalarConditions}, with flux
    entries and flows), and the thermo module keeps a third, private set
    on its temperature space.  One condition set over the fields of a
    problem would replace the three.
  \item The physics owns its space, essential dofs and loads.  A coupling
    layer needs these behind the common interface (\code{QuasiStaticProblem}:
    residual, Jacobian, pseudo-time, Dirichlet application) plus field
    exchange by name through the field registry.
  \item Time integration is a decorator over that interface, not a term of
    the kernels: with the unknown of the new time level as the variable,
    a second-order system adds $c_M\bM$ and a history vector to the static
    residual (Section~\ref{sec:dynamics}).  The scalar transport, whose
    capacity may depend on the state, integrates its rate term inside the
    kernel instead, by implicit Euler from an accepted grid function; a
    first-order decorator with the generalized-$\alpha$ parameters and the
    same $\rho_\infty$ as the solid's, which is what a monolithic coupling
    of the two needs, would apply to its constant-capacity case.
  \item Materials are stateless, one model per problem with region-wise
    parameters; internal variables, temperature dependence and hand-coded
    tangents are absent by design.  The kinematics (finite or small strain)
    is a trait of the material, not of the kernel: the flux contract is
    shared, and what depends on the kinematics outside the flux goes through
    \file{materials/kinematics.hpp}.  An eigenstrain (thermal, shrinkage)
    would enter where the small-strain material forms $\beps$, once
    materials can read a field.
  \item All kernels are CPU host code written without allocation or
    virtual calls in the quadrature loops, so device annotations and
    \code{mfem::forall} can be added without restructuring; partial
    assembly is not implemented.
\end{itemize}

%============================================================================
\appendix
\section{Symbols and input keys}
%============================================================================

\begin{center}
\small
\begin{tabular}{@{}l p{6.2cm} p{4.6cm}@{}}
\toprule
Symbol & Meaning & Input key \\
\midrule
$k$ & polynomial order of the displacement space & \code{mesh.order} \\
$E, \nu$ & Young's modulus, Poisson ratio & \code{material.E}, \code{material.nu} \\
$\mu, \kappa$ & shear and bulk modulus & \code{material.mu}, \code{material.kappa} \\
$u(J)$ & volumetric law $U = \kappa u$ & \code{material.volumetric} \\
$c_1, c_2$ & Mooney--Rivlin constants & \code{material.c1}, \code{material.c2} \\
$c_{10}, c_{20}, c_{30}$ & Yeoh constants & \code{material.c10}, \code{c20}, \code{c30} \\
$J_m$ & Gent limit of $\Ib - 3$ & \code{material.Jm} \\
$N$ & Arruda--Boyce links per chain & \code{material.N} \\
$\mathcal{L}^{-1}$ & inverse Langevin approximation, \code{pade} (default) or \code{series} & \code{material.inverse\_langevin} \\
$\mu_r, \alpha_r$ & Ogden terms & \code{material.mu\_r}, \code{material.alpha\_r} \\
$\rho_R$ & reference density, also by region & \code{material.rho0} \\
$G_i, \tau_i$ & Maxwell branches: non-equilibrium shear modulus, relaxation time & \code{material.branches} \\
$\theta_0, \alpha, c_v, k$ & reference temperature, expansion, heat capacity, conductivity & \code{material.thermal} \\
$s(\theta)$ & entropic scaling of the shear modulus & \code{material.thermal.entropic} \\
$\bar\theta, h$ & prescribed temperature, inward heat flux & \code{bcs.temperature}, \code{bcs.heat\_flux} \\
$r, z$ & coordinates of an axisymmetric mesh & \code{plane: axisymmetric} \\
$\bX_{\mathrm{pin}}$ & a pinned node & \code{point} of a Dirichlet entry \\
$t_{\text{final}}, \Delta t$ (rate-dependent, no inertia) & end time, time step & \code{time.t\_final}, \code{dt} | \code{steps} \\
$r, \bm c(t), k$ & rigid sphere: radius, centre, penalty & \code{bcs.contact} \\
$t_{\text{final}}, \Delta t$ & end time, time step & \code{dynamics.t\_final}, \code{dt} | \code{steps} \\
$\beta, \gamma$; $\alpha$; $\rho_\infty$ & time integration parameters & \code{dynamics.scheme} and its keys \\
$\bu_0, \bm v_0$ & initial displacement, velocity & \code{dynamics.initial} \\
$\bar\bu, \bar\bT, p, \bb$ & prescribed data & \code{bcs.dirichlet}, \code{bcs.traction}, \code{body\_force} \\
$s_i(t)$ & load schedule & \code{schedule} of an entry \\
$c(u), \kappa(u), s$ & capacity, conductivity (affine laws), reaction coefficient of the scalar transport & \code{transport.capacity}, \code{conductivity}, \code{reaction} \\
$\bm\beta, f$ & velocity and source of the scalar transport & \code{transport.velocity}, \code{transport.source} \\
$u_0, \bar u, g, u_{\mathrm{ex}}$ & initial condition, prescribed value, inward flux, exact solution of the scalar transport & \code{initial}, \code{bcs.dirichlet}, \code{bcs.flux}, \code{output.exact} \\
$t_1, \dots, t_n$ & load-step breakpoints & \code{solver.load\_steps} or \code{solver.steps} \\
$x^{(0)}$ & initial guess of an increment & \code{solver.predictor} \\
rtol, atol, $c$ & Newton tolerances, Armijo constant & \code{solver.newton} \\
$\gamma/\mu$ & augmentation parameter & \code{solver.linear.augmentation} \\
\bottomrule
\end{tabular}
\end{center}

\section{Source files}

\begin{center}
\small
\begin{tabular}{@{}l p{6.9cm}@{}}
\toprule
File & Content \\
\midrule
\file{base/tensor.hpp} & \code{tensor<T, n...>} of rank 1, 2, 4; products, contractions, transpose, trace, deviator, determinant and inverse up to $3\times3$ \\
\file{base/dual.hpp} & forward-mode dual number and elementary functions \\
\file{base/config.\{hpp,cpp\}} & YAML schema of the solid, validation with key paths, schedules \\
\file{base/scalar\_config.\{hpp,cpp\}}, \file{node\_reader.hpp} & YAML schema of the scalar transport; the map reader the two schemas share \\
\file{base/expression.\{hpp,cpp\}} & parser and evaluator of $f(x, y, z, t)$ \\
\file{base/coefficients.hpp} & expression and affine MFEM coefficients \\
\file{base/mesh\_input.\{hpp,cpp\}} & mesh loading, attribute resolution, jitter, refinement \\
\file{base/fields.hpp}, \file{output}, \file{probes} & field registry, ParaView writer, point probes \\
\file{kernels/total\_lagrangian.hpp} & quadrature-point functions and the displacement integrator \\
\file{kernels/mixed\_total\_lagrangian.hpp} & $u$--$p$ quadrature-point functions and block integrator \\
\file{kernels/thermo\_mixed\_total\_lagrangian.hpp} & $u$--$p$--$\theta$ kernel: point densities, dual-seeded tangent blocks, the heat-flux face integrator, the adapter of the two-block face integrators \\
\file{kernels/scalar\_flux.hpp} & scalar CG kernel: the point densities $r$ and $\bm{\mathcal F}$, the dual-seeded tangent \\
\file{kernels/follower\_pressure.hpp} & follower-pressure face integrators (both formulations) \\
\file{kernels/rigid\_sphere\_contact.hpp} & rigid-sphere penalty contact face integrators (both formulations) \\
\file{kernels/history\_field.\{hpp,cpp\}}, \file{history\_bound.hpp} & quadrature-point history of a material; its per-point binding to the kernels' stateless contract \\
\file{kernels/materials/*.hpp} & material functors (hyperelastic, \file{linear\_elastic.hpp}, the \file{viscoelastic.hpp} adapter with Maxwell branches, the \file{thermoelastic.hpp} adapter with a temperature), isochoric helpers, volumetric laws, spectral solver, plane-stress adapter, dual tangent \\
\file{kernels/materials/kinematics.hpp} & finite or small strain as a trait of the material: Cauchy stress, volume ratio, strain \\
\file{kernels/materials/scalar\_transport\_model.hpp} & the affine laws and the convection form of the scalar transport \\
\file{kernels/materials/materials.\{hpp,cpp\}} & material variants, factory, moduli resolution, region tables \\
\file{physics/loads.\{hpp,cpp\}} & load set: Dirichlet entries on faces or single nodes, dead loads, follower and contact scales, the axisymmetric weight, schedules, reactions and resultants \\
\file{physics/solid\_mechanics\_tl.\{hpp,cpp\}} & displacement formulation module \\
\file{physics/mixed\_solid\_mechanics\_tl.\{hpp,cpp\}} & mixed formulation module \\
\file{physics/thermo\_solid\_mechanics\_tl.\{hpp,cpp\}} & coupled $u$--$p$--$\theta$ module: the third space, the accepted $\bC$ and $\theta$, thermal entries, the initial temperature \\
\file{physics/solid\_problem.\{hpp,cpp\}} & common interface, factory by formulation, YAML load installer \\
\file{physics/scalar\_transport.\{hpp,cpp\}} & scalar transport module: the accepted state and the step, linearity and the reuse of the operator, errors, flows, fields \\
\file{physics/scalar\_conditions.\{hpp,cpp\}} & Dirichlet and flux entries of a scalar unknown, the flows \\
\file{physics/quadrature\_fields.\{hpp,cpp\}} & quadrature quantities and their presentations \\
\file{physics/dynamic\_solid\_problem.\{hpp,cpp\}} & inertia as a decorator over a solid problem: mass matrix, step equation, initial state, energies, reactions with inertia \\
\file{solvers/time\_integration.\{hpp,cpp\}} & parameters of the generalized-$\alpha$ family, spectral radius at infinity \\
\file{solvers/quasi\_static.\{hpp,cpp\}} & load stepping and time stepping (quasi-static in physical time, dynamic) with bisection \\
\file{solvers/newton.\{hpp,cpp\}} & damped Newton with Armijo backtracking; acceptance of a linear problem at the floor of its residual \\
\file{solvers/linear\_solver.\{hpp,cpp\}} & GMRES/CG + BoomerAMG with the systems fallback \\
\file{solvers/saddle\_point\_solver.\{hpp,cpp\}} & augmented-Lagrangian FGMRES for the $u$--$p$ Jacobian; inertial part of the Schur complement approximation \\
\bottomrule
\end{tabular}
\end{center}

%============================================================================
\small\begin{thebibliography}{99}
%============================================================================
\bibitem{bonetwood} J. Bonet, R.\,D. Wood, \emph{Nonlinear Continuum Mechanics for Finite Element Analysis}, 2nd ed., Cambridge University Press, 2008.
\bibitem{holzapfel} G.\,A. Holzapfel, \emph{Nonlinear Solid Mechanics: A Continuum Approach for Engineering}, Wiley, 2000.
\bibitem{ogden} R.\,W. Ogden, \emph{Non-Linear Elastic Deformations}, Ellis Horwood, 1984 (Dover, 1997).
\bibitem{simotaylor} J.\,C. Simo, R.\,L. Taylor, Quasi-incompressible finite elasticity in principal stretches: continuum basis and numerical algorithms, \emph{Comput. Methods Appl. Mech. Engrg.} 85 (1991) 273--310.
\bibitem{pencegou} T.\,J. Pence, K. Gou, On compressible versions of the incompressible neo-Hookean material, \emph{Math. Mech. Solids} 20 (2015) 157--182.
\bibitem{doll} S. Doll, K. Schweizerhof, On the development of volumetric strain energy functions, \emph{J. Appl. Mech.} 67 (2000) 17--21.
\bibitem{ab} E.\,M. Arruda, M.\,C. Boyce, A three-dimensional constitutive model for the large stretch behavior of rubber elastic materials, \emph{J. Mech. Phys. Solids} 41 (1993) 389--412.
\bibitem{gent} A.\,N. Gent, A new constitutive relation for rubber, \emph{Rubber Chem. Technol.} 69 (1996) 59--61.
\bibitem{yeoh} O.\,H. Yeoh, Some forms of the strain energy function for rubber, \emph{Rubber Chem. Technol.} 66 (1993) 754--771.
\bibitem{herrmann} L.\,R. Herrmann, Elasticity equations for incompressible and nearly incompressible materials by a variational theorem, \emph{AIAA J.} 3 (1965) 1896--1900.
\bibitem{benzi} M. Benzi, M.\,A. Olshanskii, An augmented Lagrangian-based approach to the Oseen problem, \emph{SIAM J. Sci. Comput.} 28 (2006) 2095--2113.
\bibitem{cohen} A. Cohen, A Pad\'e approximant to the inverse Langevin function, \emph{Rheol. Acta} 30 (1991) 270--273.
\bibitem{newmark} N.\,M. Newmark, A method of computation for structural dynamics, \emph{J. Eng. Mech. Div. ASCE} 85 (1959) 67--94.
\bibitem{hht} H.\,M. Hilber, T.\,J.\,R. Hughes, R.\,L. Taylor, Improved numerical dissipation for time integration algorithms in structural dynamics, \emph{Earthquake Eng. Struct. Dyn.} 5 (1977) 283--292.
\bibitem{chunghulbert} J. Chung, G.\,M. Hulbert, A time integration algorithm for structural dynamics with improved numerical dissipation: the generalized-$\alpha$ method, \emph{J. Appl. Mech.} 60 (1993) 371--375.
\bibitem{cahouetchabard} J. Cahouet, J.-P. Chabard, Some fast 3D finite element solvers for the generalized Stokes problem, \emph{Int. J. Numer. Methods Fluids} 8 (1988) 869--895.
\bibitem{mfem} R. Anderson et al., MFEM: A modular finite element methods library, \emph{Comput. Math. Appl.} 81 (2021) 42--74.
\bibitem{anandbook} L. Anand, \emph{Introduction to Coupled Theories in Solid Mechanics}, Oxford University Press, 2025; companion codes by E.\,M. Stewart, S.\,A. Chester and L. Anand, \url{solidmechanicscoupledtheories.github.io}.
\bibitem{wang2016} S. Wang, M. Decker, D.\,L. Henann, S.\,A. Chester, Modeling of dielectric viscoelastomers with application to electromechanical instabilities, \emph{J. Mech. Phys. Solids} 95 (2016) 213--229.
\end{thebibliography}

\end{document}
